% generated by scripts/experiment-report.ts --all: every experiment report in one document - do not edit
\documentclass{polyexp}
\titleformat{\subsection}{\large\bfseries\sffamily}{}{0pt}{}
\titleformat{\subsubsection}{\normalsize\bfseries\sffamily\color{expgrey}}{}{0pt}{}
\setcounter{tocdepth}{1}
\begin{document}
{\sffamily\Huge\bfseries Polytopia AI research}\par\medskip
{\sffamily\Large Experiment notebook}\par\medskip
{\sffamily\small\color{expgrey}12 experiments \textperiodcentered\ built 2026-11-25}\par\bigskip
\tableofcontents
\clearpage\phantomsection\addcontentsline{toc}{section}{Experiment 1: Defend threatened cities (tie)}

% generated by scripts/experiment-report.ts from experiment.json and models/<run>/ - do not edit
\providecommand{\ExpHeader}{}\renewcommand{\ExpHeader}{%
{\sffamily\footnotesize\color{expgrey}EXPERIMENT 1 \textperiodcentered\ MECHANISM \hfill 2026-11-24}\par
{\LARGE\bfseries\sffamily Defend threatened cities}\hfill\verdictbox{expgrey}{TIE, NOT KEPT}\par\smallskip
{\itshape Does keeping a unit in a threatened city make the bot stronger?}\par\medskip
\begin{tabular}{@{}p{22mm}p{\dimexpr\linewidth-26mm}@{}}
\fact{Change}{A move into an empty threatened own city scores +defend; a move that empties one scores $-$defend. A city is threatened when an enemy unit is within defendRadius tiles.}
\fact{Knobs}{\knob{defend} $\in$ \{30\}, best 37.9; \knob{defendRadius} $\in$ \{2\}, best 1}
\fact{Base}{\model{greedy-full-v1}}
\fact{Method}{Tuning run (hill climbing), 100,000 games, against a control run of the same budget; metric: points share}
\fact{Setups}{4 players, 16$\times$16, drylands, turn cap 150}
\fact{Record}{issue \issue{28}; status: done}
\end{tabular}\par\medskip
\begin{center}\begin{tikzpicture}[node distance=5mm]
    \node[flowbox] (n0) {{\sffamily\bfseries\scriptsize\color{expgrey}BASE}\\[1pt]\model{full-v1}};
    \node[flowbox,right=5mm of n0] (n1) {{\sffamily\bfseries\scriptsize\color{expgrey}CHANGE}\\[1pt]\knob{defend}\\{\scriptsize 30}\\\knob{defendRadius}\\{\scriptsize 2}};
    \node[flowbox,right=5mm of n1] (n2) {{\sffamily\bfseries\scriptsize\color{expgrey}MEASURED}\\[1pt]Tuning run, 100,000 games, vs a control\\{\scriptsize 1 setup, metric: points share}};
    \node[flowbox,right=5mm of n2] (n3) {{\sffamily\bfseries\scriptsize\color{expgrey}RESULT}\\[1pt]Points vs base, final test\\{\scriptsize 50.0 $\rightarrow$ 50.7\,\%}\\\textbf{+0.7\pp}};
    \node[flowbox,right=5mm of n3] (n4) {{\sffamily\bfseries\scriptsize\color{expgrey}VERDICT}\\[1pt]\textbf{\color{expgrey}TIE, NOT KEPT}};
    \draw[flowarrow] (n0) -- (n1); \draw[flowarrow] (n1) -- (n2); \draw[flowarrow] (n2) -- (n3); \draw[flowarrow] (n3) -- (n4);
\end{tikzpicture}\end{center}
\begin{center}\sffamily\small\begin{tabular}{@{}lrrr@{}}\toprule
Measure & Base & New & Change\\\midrule
Points vs base, final test & 50.0\,\% & 50.7\,\% & +0.7\\
Won vs base, final test & 50.0\,\% & 51.0\,\% & +1.0\\
Points share, 6-model panel & 67.7\,\% & 67.6\,\% & $-$0.1\\
\bottomrule\end{tabular}\end{center}
}
\providecommand{\ExpResults}{}\renewcommand{\ExpResults}{%
\subsection{Tuning run, 100,000 games {\normalfont\small(\model{greedy-full-defend})}}
\begin{center}\begin{tikzpicture}\begin{axis}[expplot,height=4.6cm,xlabel={training games (thousands)},ylabel={Points share, panel mean (\%)},enlarge y limits=0.2]
\addplot[dashed,expgrey,line width=0.7pt,forget plot] coordinates {(0.0,65.90) (100.0,65.90)};
\addplot[expblue,mark=*,mark size=1.8pt] coordinates {(0.0,65.90) (12.9,67.40) (25.1,67.40) (37.7,67.40) (50.4,67.40) (62.9,67.40) (75.2,67.40) (87.5,67.60) (100.0,67.60)};
\end{axis}\end{tikzpicture}\\[-2pt]{\footnotesize\color{expgrey}Learning curve of the tuning run (3 kept changes in 320 rounds). Dashed: the start, i.e. the mechanism at its hand-set value.}\end{center}
\begin{center}\sffamily\small\begin{tabular}{@{}lrr@{}}\toprule
Opponent & Won (\%) & Points (\%)\\\midrule
\model{v1} & 95 $\pm$ 2 & 75.2\\
\model{hill-v1} & 86 $\pm$ 3 & 69.1\\
\model{hill-v2} & 94 $\pm$ 2 & 71.4\\
\model{hill-v3} & 87 $\pm$ 3 & 67.7\\
\model{hill-v4} & 88 $\pm$ 3 & 67.2\\
\model{full-v1} & 51 $\pm$ 5 & 50.7\\
\midrule
\textbf{Mean} & 83.3 & 66.9\\
\bottomrule\end{tabular}\\[3pt]{\footnotesize\color{expgrey}Final test of the saved model: 400 unseen games against each panel model.}\end{center}
}

\ExpHeader

\subsection{Question and hypothesis}
No part of the full-action bot keeps a unit in a city. Every unit is pulled toward the nearest village, ruin, enemy city or enemy unit, so cities next to enemies are often empty and captured for free. \textbf{Hypothesis:} a garrison rule stops these free captures and raises strength against \model{greedy-full-v1}.

\subsection{Mechanism}
A city is \emph{threatened} when an enemy unit (not at peace) is within \knob{defendRadius} tiles. A move into an empty threatened own city scores $+\knob{defend}$; a move that leaves one empty scores $-\knob{defend}$, unless a capture or attack elsewhere outweighs it. With \knob{defend} $=0$ the bot plays exactly like the base. Code: \texttt{engine/core/src/bot\_full.rs}.

\subsection{Setup}
First experiment of the series (\issue{31}), run under the original protocol: one hill-climbing run of 100,000 games from the base (seed 4242, league \texttt{self,past}, batch 30), compared with a control run of the same budget and no new mechanism (\model{greedy-full-v2-ctrl}: 66.0\,\% on the panel at 100,000 games). The knobs started at hand-set values: \knob{defend} 30, \knob{defendRadius} 2.

\subsection{Results}
\ExpResults

Untrained, the mechanism hurts: at the hand-set values the panel mean is 65.9\,\%, 1.8 points below the base (67.7\,\%), and 46.1\,\% of points against the base. Training repaired it. The first kept change cut the radius from 2 to 1; later kept changes raised \knob{defend} from 30 to 37.9. The trained model reaches 67.6\,\% on the panel and 50.7\,\% of points against the base (51\,\% won, $\pm$5), which is 1.6 points above the control and level with the base.

\subsection{Behaviour}
The behaviour metric planned for this step was not measured; the strength result was judged clear enough on its own. This is a gap in the record. Planned metric: cities lost while empty.

\subsection{Interpretation}
Holding units back for every enemy within two tiles costs more tempo than the saved cities are worth. Training wants the rule only against enemies directly next to the city, where it is neutral. The hypothesis (free captures are a major loss) is not supported at this strength level: either they are rarer than assumed, or a unit that stays home is worth about as much as the city it saves.

\todo{jakob}{Your reading: does this match what you saw? Anything to add or correct?}

\subsection{Threats to validity}
\begin{itemize}
\item One run with one seed; hill climbing kept only 3 changes in 320 rounds, so the tuned values are weakly determined.
\item Measured before the finish score existed; the effect on game length and conquest rate is unknown.
\item The behaviour metric was not measured, so the claim about free captures rests on the strength result alone.
\end{itemize}

\subsection{Verdict and next step}
\textbf{Tie, not kept.} The code stays, switched off by default. The champion remains \model{greedy-full-v1}. Next: research usefulness (experiment 2).

\subsection{Notes for the paper}
Draft: a plausible hand-written heuristic (garrison threatened cities) is neutral once tuned and harmful at its hand-set size. Supports the general point that mechanism strength must be measured, not assumed.

\todo{jakob}{Is this worth a place in the paper? Where?}


\clearpage\phantomsection\addcontentsline{toc}{section}{Experiment 2: Score research by usefulness (not kept)}

% generated by scripts/experiment-report.ts from experiment.json and models/<run>/ - do not edit
\providecommand{\ExpHeader}{}\renewcommand{\ExpHeader}{%
{\sffamily\footnotesize\color{expgrey}EXPERIMENT 2 \textperiodcentered\ MECHANISM \hfill 2026-11-24}\par
{\LARGE\bfseries\sffamily Score research by usefulness}\hfill\verdictbox{expred}{NOT KEPT}\par\smallskip
{\itshape Does the bot get stronger when it prefers techs it can use right away?}\par\medskip
\begin{tabular}{@{}p{22mm}p{\dimexpr\linewidth-26mm}@{}}
\fact{Change}{The research score gains techUse per action the tech would unlock right now, and loses techUseless when it unlocks nothing.}
\fact{Knobs}{\knob{techUse} $\in$ \{3\}, best 3; \knob{techUseless} $\in$ \{15\}, best 12.8}
\fact{Base}{\model{greedy-full-v1}}
\fact{Method}{Tuning run (hill climbing), 100,000 games, against a control run of the same budget; metric: points share}
\fact{Setups}{4 players, 16$\times$16, drylands, turn cap 150}
\fact{Record}{issue \issue{29}, pull request \issue{36}; status: done}
\end{tabular}\par\medskip
\begin{center}\begin{tikzpicture}[node distance=5mm]
    \node[flowbox] (n0) {{\sffamily\bfseries\scriptsize\color{expgrey}BASE}\\[1pt]\model{full-v1}};
    \node[flowbox,right=5mm of n0] (n1) {{\sffamily\bfseries\scriptsize\color{expgrey}CHANGE}\\[1pt]\knob{techUse}\\{\scriptsize 3}\\\knob{techUseless}\\{\scriptsize 15}};
    \node[flowbox,right=5mm of n1] (n2) {{\sffamily\bfseries\scriptsize\color{expgrey}MEASURED}\\[1pt]Tuning run, 100,000 games, vs a control\\{\scriptsize 1 setup, metric: points share}};
    \node[flowbox,right=5mm of n2] (n3) {{\sffamily\bfseries\scriptsize\color{expgrey}RESULT}\\[1pt]Points vs base, final test\\{\scriptsize 50.0 $\rightarrow$ 39.8\,\%}\\\textbf{$-$10.2\pp}};
    \node[flowbox,right=5mm of n3] (n4) {{\sffamily\bfseries\scriptsize\color{expgrey}VERDICT}\\[1pt]\textbf{\color{expred}NOT KEPT}};
    \draw[flowarrow] (n0) -- (n1); \draw[flowarrow] (n1) -- (n2); \draw[flowarrow] (n2) -- (n3); \draw[flowarrow] (n3) -- (n4);
\end{tikzpicture}\end{center}
\begin{center}\sffamily\small\begin{tabular}{@{}lrrr@{}}\toprule
Measure & Base & New & Change\\\midrule
Points vs base, final test & 50.0\,\% & 39.8\,\% & $-$10.2\\
Won vs base, final test & 50.0\,\% & 28.2\,\% & $-$21.8\\
Points share, 6-model panel & 67.7\,\% & 63.3\,\% & $-$4.4\\
\bottomrule\end{tabular}\end{center}
}
\providecommand{\ExpResults}{}\renewcommand{\ExpResults}{%
\subsection{Tuning run, 100,000 games {\normalfont\small(\model{greedy-full-research})}}
\begin{center}\begin{tikzpicture}\begin{axis}[expplot,height=4.6cm,xlabel={training games (thousands)},ylabel={Points share, panel mean (\%)},enlarge y limits=0.2]
\addplot[dashed,expgrey,line width=0.7pt,forget plot] coordinates {(0.0,63.40) (100.0,63.40)};
\addplot[expblue,mark=*,mark size=1.8pt] coordinates {(0.0,63.40) (12.9,63.30) (25.2,63.30) (37.7,63.30) (50.3,63.30) (62.9,63.30) (75.2,63.30) (87.7,63.30) (100.0,63.30)};
\end{axis}\end{tikzpicture}\\[-2pt]{\footnotesize\color{expgrey}Learning curve of the tuning run (1 kept changes in 244 rounds). Dashed: the start, i.e. the mechanism at its hand-set value.}\end{center}
\begin{center}\sffamily\small\begin{tabular}{@{}lrr@{}}\toprule
Opponent & Won (\%) & Points (\%)\\\midrule
\model{v1} & 84 $\pm$ 4 & 69.6\\
\model{hill-v1} & 79 $\pm$ 4 & 65.0\\
\model{hill-v2} & 81 $\pm$ 4 & 65.0\\
\model{hill-v3} & 76 $\pm$ 4 & 60.8\\
\model{hill-v4} & 76 $\pm$ 4 & 60.8\\
\model{full-v1} & 28 $\pm$ 4 & 39.8\\
\midrule
\textbf{Mean} & 70.8 & 60.2\\
\bottomrule\end{tabular}\\[3pt]{\footnotesize\color{expgrey}Final test of the saved model: 400 unseen games against each panel model.}\end{center}
}

\ExpHeader

\subsection{Question and hypothesis}
Research is scored as a fixed weight minus cost and tier. It never asks whether the tech unlocks anything the player can use, so bots buy Fishing with no water in reach. \textbf{Hypothesis:} scoring a tech by what it would let the player do now avoids wasted stars and raises strength.

\subsection{Mechanism}
$\mathit{uses}(t)$ counts the actions tech $t$ would add right now: builds and harvests on own tiles it unlocks, unit types a city could train, abilities with a valid target. The research score gets $+\,\knob{techUse}\cdot\mathit{uses}(t) - \knob{techUseless}\cdot[\mathit{uses}(t)=0]$. Both knobs default to 0.

\subsection{Setup}
Same protocol as experiment 1: one hill-climbing run of 100,000 games from the base against the same control. Hand-set start: \knob{techUse} 3, \knob{techUseless} 15. The research score these compete with is about 28 minus cost and tier, so the start values were large against it.

\subsection{Results}
\ExpResults

The mechanism costs 4.3 points on the panel at its hand-set values (63.4\,\% against 67.7\,\%) and loses clearly to the base: 39.8\,\% of points, 28\,\% won ($\pm$4). Training could not repair it: one kept change in 244 rounds (\knob{techUseless} 15 to 12.8).

\subsection{Behaviour}
The behaviour metric planned for this step was not measured; the strength result was judged clear enough on its own. This is a gap in the record. Planned metric: share of researched techs whose unlocks are never used.

\subsection{Interpretation}
``What can I build right now'' undervalues techs whose payoff comes later (a unit tech before a war, the parent of a needed tech) and overvalues cheap terrain builds. The base had already tuned its research order through per-tech weights, so the new term mostly overrode a tuned order. The start was also about five times too strong, and multiplicative nudges of $\pm$40\,\% cannot walk a knob that far within the budget. This experiment is the reason later sweeps always include a weak value.

\todo{jakob}{Your reading: does this match what you saw? Anything to add or correct?}

\subsection{Threats to validity}
\begin{itemize}
\item Only one strength was effectively tested. A mild version (for example \knob{techUse} 0.3, \knob{techUseless} 2) is untested and could behave differently.
\item The behaviour metric was not measured.
\end{itemize}

\subsection{Verdict and next step}
\textbf{Not kept.} The code stays, switched off. A retry at a much weaker strength would be a new experiment. Next: front pull (experiment 3).

\subsection{Notes for the paper}
Draft: a myopic usefulness heuristic for research loses 10 points of share against tuned per-tech weights; and hill climbing cannot recover from a badly sized start. Motivates the move from tuning runs to strength sweeps.

\todo{jakob}{Is this worth a place in the paper? Where?}


\clearpage\phantomsection\addcontentsline{toc}{section}{Experiment 3: Pull idle units to the front (kept)}

% generated by scripts/experiment-report.ts from experiment.json and models/<run>/ - do not edit
\providecommand{\ExpHeader}{}\renewcommand{\ExpHeader}{%
{\sffamily\footnotesize\color{expgrey}EXPERIMENT 3 \textperiodcentered\ MECHANISM \hfill 2026-11-24}\par
{\LARGE\bfseries\sffamily Pull idle units to the front}\hfill\verdictbox{expgreen}{KEPT}\par\smallskip
{\itshape Does the bot get stronger when idle units drift toward the nearest enemy?}\par\medskip
\begin{tabular}{@{}p{22mm}p{\dimexpr\linewidth-26mm}@{}}
\fact{Change}{A move scores +front per tile it brings the unit closer to the nearest known enemy city or unit. The centre pull remains the fallback when no enemy is known.}
\fact{Knobs}{\knob{front} $\in$ \{0.5, 1, 2, 4, 8\}, chosen 1.13}
\fact{Base}{\model{greedy-full-v1} $\rightarrow$ result \model{greedy-full-front}}
\fact{Method}{Tuning run (100,000 games), then the first strength sweep (200 games per point, final test 400); metric: points share}
\fact{Setups}{4 players, 16$\times$16, drylands, turn cap 150}
\fact{Record}{issue \issue{30}, pull request \issue{62}; status: done}
\end{tabular}\par\medskip
\begin{center}\begin{tikzpicture}[node distance=5mm]
    \node[flowbox] (n0) {{\sffamily\bfseries\scriptsize\color{expgrey}BASE}\\[1pt]\model{full-v1}};
    \node[flowbox,right=5mm of n0] (n1) {{\sffamily\bfseries\scriptsize\color{expgrey}CHANGE}\\[1pt]\knob{front}\\{\scriptsize 0.5, 1, 2, 4, 8}};
    \node[flowbox,right=5mm of n1] (n2) {{\sffamily\bfseries\scriptsize\color{expgrey}MEASURED}\\[1pt]Tuning run, then the first strength sweep\\{\scriptsize 1 setup, metric: points share}};
    \node[flowbox,right=5mm of n2] (n3) {{\sffamily\bfseries\scriptsize\color{expgrey}RESULT}\\[1pt]Points vs base, final test\\{\scriptsize 50.0 $\rightarrow$ 54.8\,\%}\\\textbf{+4.8\pp}};
    \node[flowbox,right=5mm of n3] (n4) {{\sffamily\bfseries\scriptsize\color{expgrey}VERDICT}\\[1pt]\textbf{\color{expgreen}KEPT}\\{\scriptsize new champion:}\\\model{full-front}};
    \draw[flowarrow] (n0) -- (n1); \draw[flowarrow] (n1) -- (n2); \draw[flowarrow] (n2) -- (n3); \draw[flowarrow] (n3) -- (n4);
\end{tikzpicture}\end{center}
\begin{center}\sffamily\small\begin{tabular}{@{}lrrr@{}}\toprule
Measure & Base & New & Change\\\midrule
Points vs base, final test & 50.0\,\% & 54.8\,\% & +4.8\\
Won vs base, final test & 50.0\,\% & 57.8\,\% & +7.8\\
Points share, 6-model panel & 67.7\,\% & 70.3\,\% & +2.6\\
\bottomrule\end{tabular}\end{center}
}
\providecommand{\ExpResults}{}\renewcommand{\ExpResults}{%
\subsection{Tuning run, 100,000 games (the new champion) {\normalfont\small(\model{greedy-full-front})}}
\begin{center}\begin{tikzpicture}\begin{axis}[expplot,height=4.6cm,xlabel={training games (thousands)},ylabel={Points share, panel mean (\%)},enlarge y limits=0.2]
\addplot[dashed,expgrey,line width=0.7pt,forget plot] coordinates {(0.0,68.10) (100.3,68.10)};
\addplot[expblue,mark=*,mark size=1.8pt] coordinates {(0.0,68.10) (12.5,68.10) (25.4,68.10) (37.6,68.10) (50.2,68.10) (62.9,68.10) (75.4,70.30) (87.8,70.30) (100.3,70.30)};
\end{axis}\end{tikzpicture}\\[-2pt]{\footnotesize\color{expgrey}Learning curve of the tuning run (2 kept changes in 258 rounds). Dashed: the start, i.e. the mechanism at its hand-set value.}\end{center}
\begin{center}\sffamily\small\begin{tabular}{@{}lrr@{}}\toprule
Opponent & Won (\%) & Points (\%)\\\midrule
\model{v1} & 94 $\pm$ 2 & 75.7\\
\model{hill-v1} & 91 $\pm$ 3 & 72.3\\
\model{hill-v2} & 95 $\pm$ 2 & 73.7\\
\model{hill-v3} & 91 $\pm$ 3 & 70.3\\
\model{hill-v4} & 87 $\pm$ 3 & 68.6\\
\model{full-v1} & 58 $\pm$ 5 & 54.8\\
\midrule
\textbf{Mean} & 85.9 & 69.2\\
\bottomrule\end{tabular}\\[3pt]{\footnotesize\color{expgrey}Final test of the saved model: 400 unseen games against each panel model.}\end{center}

\subsection{Strength sweep, untrained {\normalfont\small(\model{greedy-full-front-sweep})}}
\begin{center}\sffamily\small\begin{tabular}{@{}lrrrrrr@{}}\toprule
Point & Finish (\%) & Points (\%) & Won (\%) & Conquest (\%) & Turns & vs base: points (\%)\\\midrule
base & -- & 67.73 & 84.3 & -- & -- & 50.0\\
front=0.5 & -- & 68.50 & 86.0 & -- & -- & 53.0\\
front=1 & -- & 68.13 & 85.3 & -- & -- & 53.3\\
front=2 & -- & 68.77 & 86.8 & -- & -- & 52.9\\
front=4 & -- & 68.43 & 86.2 & -- & -- & 52.1\\
front=8 & -- & 68.55 & 86.7 & -- & -- & 53.4\\
\bottomrule\end{tabular}\\[3pt]{\footnotesize\color{expgrey}Panel means; 200 games against each of 6 panel models per point.}\end{center}
\begin{center}\begin{tikzpicture}\begin{axis}[expplot,height=4.6cm,xtick={0,1,2,3,4,5},xticklabels={{base},{0.5},{1},{2},{4},{8}},
  xlabel={\texttt{front}},ylabel={Points share, panel mean (\%)},enlarge x limits=0.12,enlarge y limits=0.25]
\addplot[dashed,expgrey,line width=0.7pt,forget plot] coordinates {(0,67.733) (5,67.733)};
\addplot[only marks,mark=*,mark size=2.6pt,expblue] coordinates { (1,68.500) (2,68.133) (3,68.767) (4,68.433) (5,68.550)};
\addplot[only marks,mark=*,mark size=2.6pt,expgrey] coordinates {(0,67.733)};
\end{axis}\end{tikzpicture}\\[-2pt]{\footnotesize\color{expgrey}Dashed: the base. Grey: base; blue: sweep points.}\end{center}
\begin{center}\sffamily\small\begin{tabular}{@{}lrr@{}}\toprule
Opponent & Won (\%) & Points (\%)\\\midrule
\model{v1} & 96 $\pm$ 2 & 75.5\\
\model{hill-v1} & 94 $\pm$ 2 & 72.2\\
\model{hill-v2} & 95 $\pm$ 2 & 72.3\\
\model{hill-v3} & 91 $\pm$ 3 & 69.1\\
\model{hill-v4} & 88 $\pm$ 3 & 67.0\\
\model{full-v1} & 60 $\pm$ 5 & 54.1\\
\midrule
\textbf{Mean without \model{full-v1}} & 92.5 & 71.2\\
\bottomrule\end{tabular}\\[3pt]{\footnotesize\color{expgrey}Final test of front=2: 400 unseen games against each.}\end{center}
}

\ExpHeader

\subsection{Question and hypothesis}
A move scores only if it brings the unit closer to its goal. A unit whose goal is far away or unreachable stays where it is, often at the board edge. \textbf{Hypothesis:} a pull toward the nearest enemy puts idle units to work and raises strength.

\subsection{Mechanism}
Let $d(x)$ be the distance from tile $x$ to the nearest enemy city or unit the player has seen (not at peace). A move from $a$ to $b$ gets $+\,\knob{front}\cdot(d(a)-d(b))$. Default 0. Code: \texttt{enemy\_dist}, \texttt{front\_term} in \texttt{bot\_full.rs}.

\subsection{Setup}
The tuning run started from \knob{front} $=1$ under the protocol of experiments 1 and 2. After it, the strength sweep was introduced: the base plus the knob at five hand-set values, each against the panel on the same maps as the curve points, then a final test of the best point. The sweep took 1\,min 54\,s; a tuning run takes about 30 minutes.

\subsection{Results}
\ExpResults

The trained model reaches 70.3\,\% on the panel (base 67.7\,\%, control 66.0\,\%) and beats the base head to head: 54.8\,\% of points, 58\,\% won ($\pm$5). The sweep shows the gain hardly depends on the strength: every value from 0.5 to 8 adds 0.4 to 1.1 points on the panel and about 53\,\% against the base. Training added about 1.5 points more, from two kept changes late in the run that also moved \knob{open.attack}, \knob{greedy.richAt} and \knob{greedy.hit}.

\subsection{Behaviour}
The behaviour metric planned for this step was not measured; the strength result was judged clear enough on its own. This is a gap in the record. Planned metric: mean distance of units to the nearest enemy.

\subsection{Interpretation}
The pull mostly gets idle units moving at all, and any strength does that; this is why the response is flat across a factor of 16. The first mechanism of the series that adds something the bot could not do before, instead of re-weighting a decision it already made.

\todo{jakob}{Your reading: does this match what you saw? Anything to add or correct?}

\subsection{Threats to validity}
\begin{itemize}
\item The champion's extra 1.5 points come from two late kept changes in one run; a different seed might not find them.
\item Under the finish score (experiment 5) a stronger pull trades points for faster conquests, so ``flat'' holds for points share only.
\end{itemize}

\subsection{Verdict and next step}
\textbf{Kept} (head-to-head rule). New champion: \model{greedy-full-front}. Process change: new mechanisms are measured with a strength sweep; tuning runs happen every three to four kept mechanisms. Next: road-aware moves (experiment 4).

\subsection{Notes for the paper}
Draft: mechanisms that enable a new behaviour (move when idle) give gains that are insensitive to their weight; a 2-minute sweep finds them as reliably as a 30-minute tuning run. Candidate figure: the sweep plot (flat response) next to the learning curve (170 rounds with nothing kept).

\todo{jakob}{Is this worth a place in the paper? Where?}


\clearpage\phantomsection\addcontentsline{toc}{section}{Experiment 4: Road-aware moves (not kept)}

% generated by scripts/experiment-report.ts from experiment.json and models/<run>/ - do not edit
\providecommand{\ExpHeader}{}\renewcommand{\ExpHeader}{%
{\sffamily\footnotesize\color{expgrey}EXPERIMENT 4 \textperiodcentered\ MECHANISM \hfill 2026-11-24}\par
{\LARGE\bfseries\sffamily Road-aware moves}\hfill\verdictbox{expred}{NOT KEPT}\par\smallskip
{\itshape Does the bot get stronger when moves count real movement cost, so that roads are used?}\par\medskip
\begin{tabular}{@{}p{22mm}p{\dimexpr\linewidth-26mm}@{}}
\fact{Change}{A move's progress becomes tiles gained + roads $\times$ (road-aware gain $-$ tiles gained); steps between two road tiles cost half, as in the engine.}
\fact{Knobs}{\knob{roads} $\in$ \{0.25, 0.5, 1, 2\}, best 0.5}
\fact{Base}{\model{greedy-full-front}}
\fact{Method}{Strength sweep, 200 games per point against each panel model, final test 400 games; metric: points share}
\fact{Setups}{4 players, 16$\times$16, drylands, turn cap 150}
\fact{Record}{issue \issue{34}, pull request \issue{125}; status: done}
\end{tabular}\par\medskip
\begin{center}\begin{tikzpicture}[node distance=5mm]
    \node[flowbox] (n0) {{\sffamily\bfseries\scriptsize\color{expgrey}BASE}\\[1pt]\model{full-front}};
    \node[flowbox,right=5mm of n0] (n1) {{\sffamily\bfseries\scriptsize\color{expgrey}CHANGE}\\[1pt]\knob{roads}\\{\scriptsize 0.25, 0.5, 1, 2}};
    \node[flowbox,right=5mm of n1] (n2) {{\sffamily\bfseries\scriptsize\color{expgrey}MEASURED}\\[1pt]Strength sweep, 4 values\\{\scriptsize 1 setup, metric: points share}};
    \node[flowbox,right=5mm of n2] (n3) {{\sffamily\bfseries\scriptsize\color{expgrey}RESULT}\\[1pt]Points vs base, final test\\{\scriptsize 50.0 $\rightarrow$ 50.2\,\%}\\\textbf{+0.2\pp}};
    \node[flowbox,right=5mm of n3] (n4) {{\sffamily\bfseries\scriptsize\color{expgrey}VERDICT}\\[1pt]\textbf{\color{expred}NOT KEPT}};
    \draw[flowarrow] (n0) -- (n1); \draw[flowarrow] (n1) -- (n2); \draw[flowarrow] (n2) -- (n3); \draw[flowarrow] (n3) -- (n4);
\end{tikzpicture}\end{center}
\begin{center}\sffamily\small\begin{tabular}{@{}lrrr@{}}\toprule
Measure & Base & New & Change\\\midrule
Points vs base, final test & 50.0\,\% & 50.2\,\% & +0.2\\
Won vs base, final test & 50.0\,\% & 50.7\,\% & +0.7\\
Points share, 7-model panel & 67.4\,\% & 67.4\,\% & +0.0\\
\bottomrule\end{tabular}\end{center}
}
\providecommand{\ExpResults}{}\renewcommand{\ExpResults}{%
\subsection{Strength sweep {\normalfont\small(\model{greedy-full-roads})}}
\begin{center}\sffamily\small\begin{tabular}{@{}lrrrrrr@{}}\toprule
Point & Finish (\%) & Points (\%) & Won (\%) & Conquest (\%) & Turns & vs base: points (\%)\\\midrule
base & -- & 67.36 & 80.7 & -- & -- & 50.0\\
roads=0.25 & -- & 67.37 & 80.9 & -- & -- & 49.9\\
roads=0.5 & -- & 67.40 & 80.7 & -- & -- & 50.4\\
roads=1 & -- & 67.13 & 80.1 & -- & -- & 50.2\\
roads=2 & -- & 66.74 & 79.9 & -- & -- & 50.0\\
\bottomrule\end{tabular}\\[3pt]{\footnotesize\color{expgrey}Panel means; 200 games against each of 7 panel models per point.}\end{center}
\begin{center}\begin{tikzpicture}\begin{axis}[expplot,height=4.6cm,xtick={0,1,2,3,4},xticklabels={{base},{0.25},{0.5},{1},{2}},
  xlabel={\texttt{roads}},ylabel={Points share, panel mean (\%)},enlarge x limits=0.12,enlarge y limits=0.25]
\addplot[dashed,expgrey,line width=0.7pt,forget plot] coordinates {(0,67.357) (4,67.357)};
\addplot[only marks,mark=*,mark size=2.6pt,expblue] coordinates { (1,67.371) (2,67.400) (3,67.129) (4,66.743)};
\addplot[only marks,mark=*,mark size=2.6pt,expgrey] coordinates {(0,67.357)};
\end{axis}\end{tikzpicture}\\[-2pt]{\footnotesize\color{expgrey}Dashed: the base. Grey: base; blue: sweep points.}\end{center}
\begin{center}\sffamily\small\begin{tabular}{@{}lrr@{}}\toprule
Opponent & Won (\%) & Points (\%)\\\midrule
\model{v1} & 93 $\pm$ 2 & 75.5\\
\model{hill-v1} & 91 $\pm$ 3 & 72.3\\
\model{hill-v2} & 94 $\pm$ 2 & 73.1\\
\model{hill-v3} & 90 $\pm$ 3 & 69.7\\
\model{hill-v4} & 86 $\pm$ 3 & 68.6\\
\model{full-v1} & 57 $\pm$ 5 & 54.4\\
\model{full-front} & 51 $\pm$ 5 & 50.2\\
\midrule
\textbf{Mean without \model{full-front}} & 85.1 & 68.9\\
\bottomrule\end{tabular}\\[3pt]{\footnotesize\color{expgrey}Final test of roads=0.5: 400 unseen games against each.}\end{center}
}

\ExpHeader

\subsection{Question and hypothesis}
In a live 9-bot game at turn 40 the bots had built 166 road tiles (41\,\% of the map, about 330 stars), yet only 17 of 825 moves started and ended on the road network. Move scoring counts tiles gained and ignores movement cost. \textbf{Hypothesis:} scoring moves by road-aware distance makes units use roads and reach goals sooner.

\subsection{Mechanism}
A Dijkstra map from the unit's goal gives a road-aware distance in half-steps. Progress of a move $= t + \knob{roads}\cdot(r - t)$, with $t$ the tiles gained and $r$ the road-aware gain. A step between two road tiles (road, bridge, village or city on friendly or unclaimed land) costs half. Default 0. Code: \texttt{road\_dist} in \texttt{bot\_full.rs}.

\subsection{Setup}
Sanity check on the live game at \knob{roads} $=1$: 13\,\% of units pick a different best move, and slightly more moves end on a road or city (992 against 935 of 1,247). The panel has seven models here, because the base is added to the fixed six.

\subsection{Results}
\ExpResults

No effect at weak values, a small cost at strong ones (\knob{roads} 2: $-0.6$ points). Final test of the best point against the base: 50.2\,\% of points, 51\,\% won ($\pm$5).

\subsection{Behaviour}
Moves ending on a road or city at \knob{roads} $=1$: 992 of 1,247, against 935 for the base (one live position).

\subsection{Interpretation}
Roads change which tile a unit ends on, but rarely what it achieves. On 16$\times$16 maps most moves reach their goal within one or two turns, so a half-cost path saves little. The hypothesis that unused roads are a strength problem is not supported; they are a waste of stars, and stars are not the bottleneck (2,500 banked by the late game).

\todo{jakob}{Your reading: does this match what you saw? Anything to add or correct?}

\subsection{Threats to validity}
\begin{itemize}
\item Only 16$\times$16 maps; on larger maps paths are longer and roads could matter.
\item The other half of the idea (build roads only where they connect something) is untested.
\end{itemize}

\subsection{Verdict and next step}
\textbf{Not kept.} The code stays, switched off. The champion remains \model{greedy-full-front}. Next: a yardstick that rewards finishing (experiment 5).

\subsection{Notes for the paper}
Draft: a clean negative result. Correcting a visibly ``wrong'' behaviour (roads built but unused) does not change strength, because it is not on the critical path to winning.

\todo{jakob}{Is this worth a place in the paper? Where?}


\clearpage\phantomsection\addcontentsline{toc}{section}{Experiment 5: The finish score: reward winning fast (adopted)}

% generated by scripts/experiment-report.ts from experiment.json and models/<run>/ - do not edit
\providecommand{\ExpHeader}{}\renewcommand{\ExpHeader}{%
{\sffamily\footnotesize\color{expgrey}EXPERIMENT 5 \textperiodcentered\ METHOD \hfill 2026-11-24}\par
{\LARGE\bfseries\sffamily The finish score: reward winning fast}\hfill\verdictbox{expblue}{ADOPTED}\par\smallskip
{\itshape How often do games end by conquest, and does a yardstick that rewards early wins rank bots differently from points share?}\par\medskip
\begin{tabular}{@{}p{22mm}p{\dimexpr\linewidth-26mm}@{}}
\fact{Change}{A new yardstick. A conquest win scores 75-100\% (the earlier, the more), a conquest loss 0-25\%, a game stopped at the turn cap 25-75\% by points share.}
\fact{Knobs}{\knob{front} $\in$ \{0.5, 2, 4, 8, 16, 32\}, best 8}
\fact{Base}{\model{greedy-full-front}}
\fact{Method}{Measurement of how the champion's games end, then a re-sweep of the front pull under the new metric; metric: finish score}
\fact{Setups}{4 players, 16$\times$16, drylands, turn cap 150; 8 players, 20$\times$20, drylands, turn cap 300 (self-play measurement)}
\fact{Record}{issue \issue{124}, pull request \issue{127}; status: done}
\end{tabular}\par\medskip
\begin{center}\begin{tikzpicture}[node distance=5mm]
    \node[flowbox] (n0) {{\sffamily\bfseries\scriptsize\color{expgrey}BASE}\\[1pt]\model{full-front}};
    \node[flowbox,right=5mm of n0] (n1) {{\sffamily\bfseries\scriptsize\color{expgrey}CHANGE}\\[1pt]\knob{front}\\{\scriptsize 0.5, 2, 4, 8, 16, 32}};
    \node[flowbox,right=5mm of n1] (n2) {{\sffamily\bfseries\scriptsize\color{expgrey}MEASURED}\\[1pt]How games end; front pull re-swept\\{\scriptsize 2 setups, metric: finish score}};
    \node[flowbox,right=5mm of n2] (n3) {{\sffamily\bfseries\scriptsize\color{expgrey}RESULT}\\[1pt]Finish score, 7-model panel (front 1.13 $\rightarrow$ 8)\\{\scriptsize 72.4 $\rightarrow$ 73.3\,\%}\\\textbf{+0.9\pp}};
    \node[flowbox,right=5mm of n3] (n4) {{\sffamily\bfseries\scriptsize\color{expgrey}VERDICT}\\[1pt]\textbf{\color{expblue}ADOPTED}};
    \draw[flowarrow] (n0) -- (n1); \draw[flowarrow] (n1) -- (n2); \draw[flowarrow] (n2) -- (n3); \draw[flowarrow] (n3) -- (n4);
\end{tikzpicture}\end{center}
\begin{center}\sffamily\small\begin{tabular}{@{}lrrr@{}}\toprule
Measure & Base & New & Change\\\midrule
Finish score, 7-model panel (front 1.13 $\rightarrow$ 8) & 72.4\,\% & 73.3\,\% & +0.9\\
Points vs champion, final test (front 8) & 50.0\,\% & 49.5\,\% & $-$0.5\\
Self-play games not finished by the cap, 4 players & -- & 19.0\,\% & --\\
Self-play games not finished by the cap, 8 players & -- & 36.0\,\% & --\\
\bottomrule\end{tabular}\end{center}
}
\providecommand{\ExpResults}{}\renewcommand{\ExpResults}{%
\subsection{The champion under the finish score {\normalfont\small(\model{greedy-full-front-finish})}}
\begin{center}\sffamily\small\begin{tabular}{@{}lrrrrr@{}}\toprule
Opponent & Won (\%) & Points (\%) & Finish (\%) & Conquest (\%) & Turns\\\midrule
\model{v1} & 94 $\pm$ 2 & 75.7 & 83.1 & 76 & 68\\
\model{hill-v1} & 91 $\pm$ 3 & 72.3 & 78.3 & 68 & 75\\
\model{hill-v2} & 95 $\pm$ 2 & 73.7 & 80.9 & 71 & 77\\
\model{hill-v3} & 91 $\pm$ 3 & 70.3 & 78.3 & 68 & 76\\
\model{hill-v4} & 87 $\pm$ 3 & 68.6 & 74.1 & 59 & 83\\
\model{full-v1} & 58 $\pm$ 5 & 54.8 & 55.1 & 44 & 73\\
\model{full-front} & 50 $\pm$ 5 & 50.0 & 50.0 & 40 & 69\\
\midrule
\textbf{Mean without \model{full-front}} & 85.9 & 69.2 & 75.0 & 64.0 & 75.3\\
\bottomrule\end{tabular}\\[3pt]{\footnotesize\color{expgrey}Final test of base: 400 unseen games against each.}\end{center}

\subsection{Front pull re-swept, 0.5 to 8 {\normalfont\small(\model{greedy-full-front-finish-sweep})}}
\begin{center}\sffamily\small\begin{tabular}{@{}lrrrrrr@{}}\toprule
Point & Finish (\%) & Points (\%) & Won (\%) & Conquest (\%) & Turns & vs base: points (\%)\\\midrule
base & 72.43 & 67.36 & 80.7 & 64.4 & 69.8 & 50.0\\
front=0.5 & 71.13 & 66.71 & 79.7 & 62.5 & 72.6 & 45.0\\
front=2 & 72.11 & 66.77 & 81.0 & 63.2 & 70.5 & 49.9\\
front=4 & 72.81 & 67.34 & 81.6 & 65.0 & 68.8 & 49.7\\
front=8 & 73.31 & 67.51 & 82.0 & 65.6 & 68.3 & 51.6\\
\bottomrule\end{tabular}\\[3pt]{\footnotesize\color{expgrey}Panel means; 200 games against each of 7 panel models per point.}\end{center}
\begin{center}\begin{tikzpicture}\begin{axis}[expplot,height=4.6cm,xtick={0,1,2,3,4},xticklabels={{base},{0.5},{2},{4},{8}},
  xlabel={\texttt{front}},ylabel={Finish score, panel mean (\%)},enlarge x limits=0.12,enlarge y limits=0.25]
\addplot[dashed,expgrey,line width=0.7pt,forget plot] coordinates {(0,72.429) (4,72.429)};
\addplot[only marks,mark=*,mark size=2.6pt,expblue] coordinates { (1,71.129) (2,72.114) (3,72.814) (4,73.314)};
\addplot[only marks,mark=*,mark size=2.6pt,expgrey] coordinates {(0,72.429)};
\end{axis}\end{tikzpicture}\\[-2pt]{\footnotesize\color{expgrey}Dashed: the base. Grey: base; blue: sweep points.}\end{center}
\begin{center}\sffamily\small\begin{tabular}{@{}lrrrrr@{}}\toprule
Opponent & Won (\%) & Points (\%) & Finish (\%) & Conquest (\%) & Turns\\\midrule
\model{v1} & 95 $\pm$ 2 & 76.8 & 84.1 & 77 & 65\\
\model{hill-v1} & 91 $\pm$ 3 & 73.5 & 79.8 & 72 & 71\\
\model{hill-v2} & 95 $\pm$ 2 & 73.5 & 80.0 & 69 & 77\\
\model{hill-v3} & 86 $\pm$ 3 & 68.1 & 76.5 & 67 & 75\\
\model{hill-v4} & 90 $\pm$ 3 & 68.8 & 76.0 & 63 & 80\\
\model{full-v1} & 53 $\pm$ 5 & 50.9 & 51.7 & 41 & 69\\
\model{full-front} & 51 $\pm$ 5 & 49.5 & 51.0 & 41 & 68\\
\midrule
\textbf{Mean without \model{full-front}} & 85.0 & 68.6 & 74.7 & 64.9 & 72.7\\
\bottomrule\end{tabular}\\[3pt]{\footnotesize\color{expgrey}Final test of front=8: 400 unseen games against each.}\end{center}

\subsection{Front pull re-swept, 8 to 32 {\normalfont\small(\model{greedy-full-front-finish-sweep2})}}
\begin{center}\sffamily\small\begin{tabular}{@{}lrrrrrr@{}}\toprule
Point & Finish (\%) & Points (\%) & Won (\%) & Conquest (\%) & Turns & vs base: points (\%)\\\midrule
base & 72.43 & 67.36 & 80.7 & 64.4 & 69.8 & 50.0\\
front=8 & 73.31 & 67.51 & 82.0 & 65.6 & 68.3 & 51.6\\
front=16 & 72.71 & 66.81 & 80.9 & 64.7 & 68.7 & 51.6\\
front=32 & 72.43 & 65.70 & 79.6 & 66.3 & 65.4 & 48.8\\
\bottomrule\end{tabular}\\[3pt]{\footnotesize\color{expgrey}Panel means; 200 games against each of 7 panel models per point.}\end{center}
\begin{center}\begin{tikzpicture}\begin{axis}[expplot,height=4.6cm,xtick={0,1,2,3},xticklabels={{base},{8},{16},{32}},
  xlabel={\texttt{front}},ylabel={Finish score, panel mean (\%)},enlarge x limits=0.12,enlarge y limits=0.25]
\addplot[dashed,expgrey,line width=0.7pt,forget plot] coordinates {(0,72.429) (3,72.429)};
\addplot[only marks,mark=*,mark size=2.6pt,expblue] coordinates { (1,73.314) (2,72.714) (3,72.429)};
\addplot[only marks,mark=*,mark size=2.6pt,expgrey] coordinates {(0,72.429)};
\end{axis}\end{tikzpicture}\\[-2pt]{\footnotesize\color{expgrey}Dashed: the base. Grey: base; blue: sweep points.}\end{center}
\begin{center}\sffamily\small\begin{tabular}{@{}lrrrrr@{}}\toprule
Opponent & Won (\%) & Points (\%) & Finish (\%) & Conquest (\%) & Turns\\\midrule
\model{v1} & 95 $\pm$ 2 & 76.8 & 84.1 & 77 & 65\\
\model{hill-v1} & 91 $\pm$ 3 & 73.5 & 79.8 & 72 & 71\\
\model{hill-v2} & 95 $\pm$ 2 & 73.5 & 80.0 & 69 & 77\\
\model{hill-v3} & 86 $\pm$ 3 & 68.1 & 76.5 & 67 & 75\\
\model{hill-v4} & 90 $\pm$ 3 & 68.8 & 76.0 & 63 & 80\\
\model{full-v1} & 53 $\pm$ 5 & 50.9 & 51.7 & 41 & 69\\
\model{full-front} & 51 $\pm$ 5 & 49.5 & 51.0 & 41 & 68\\
\midrule
\textbf{Mean without \model{full-front}} & 85.0 & 68.6 & 74.7 & 64.9 & 72.7\\
\bottomrule\end{tabular}\\[3pt]{\footnotesize\color{expgrey}Final test of front=8: 400 unseen games against each.}\end{center}
}

\ExpHeader

\subsection{Question and hypothesis}
Jakob, watching a live 9-bot game at turn 159: ``i feel like this is a policy stale mate''; then: ``maybe we should punish game length to win'' and ``then it might find more time efficient strategies''. The leader had 18 cities, 132 units and 2,515 stars banked, and no attack finished the game. Every training game ends at turn 150 at the latest and is judged by points share, so a conquest at turn 60 and a points lead held to turn 150 count about the same. \textbf{Hypothesis:} points share cannot see stalling; a metric that rewards early conquest will rank some settings differently.

\subsection{Mechanism}
With turn cap $T$ and a game ending at turn $t$:
\[
\text{finish} =
\begin{cases}
0.75 + 0.25\,\dfrac{T-t}{T} & \text{conquest win}\\[6pt]
0.25 - 0.25\,\dfrac{T-t}{T} & \text{conquest loss}\\[6pt]
0.25 + 0.5\cdot\text{points share} & \text{stopped at the cap}
\end{cases}
\]
Any real win beats any capped game, a fast win beats a slow one, and 50\,\% is even. Selected with \texttt{-{}-metric finish} in the trainer and the sweep; every panel result also reports conquest rate and game length.

\subsection{Setup}
The engine now returns the final turn and whether the game ended by conquest. The champion was measured in self-play in the training setup and in a setup like the live game (8 players, 20$\times$20, cap 300). Then the front pull of experiment 3 was swept again under the new metric, over a wider range.

\subsection{Results}
In self-play the champion ends 81\,\% of training-setup games by conquest (about 69 turns) and leaves 19\,\% unfinished; in the 8-player setup 64\,\% end by conquest (about 98 turns) and 36\,\% stall. Against the weaker panel models it wins 59 to 76\,\% of games by conquest.

\ExpResults

Under the finish score the front pull has an optimum near 8 (73.3\,\% against 72.4\,\% for the champion's 1.13), where points share is flat. At 32 the bot wins most often by conquest and in the fewest turns, but with the fewest points. In the final test \knob{front} $=8$ ties with the champion (finish 51.0\,\%, points 49.5\,\%, won 50\,\% $\pm$5).

\subsection{Behaviour}
Conquest rate and mean game length are now reported for every panel result; they are the behaviour metrics of this experiment. Stalling is real but partial, and worst in large games, which training never sees.

\subsection{Interpretation}
The new metric sees a trade-off that points share hides: a stronger pull converts points into faster conquests. It did not by itself produce a better bot, because the stalls have structural causes (experiments 6 and 7) that no re-weighting of existing knobs removes.

\todo{jakob}{Your reading: does this match what you saw? Anything to add or correct?}

\subsection{Threats to validity}
\begin{itemize}
\item The band limits (25\,\%, 75\,\%) are a design choice; other splits would weight speed against certainty differently.
\item A capped game scores 25 to 75\,\%, so a bot can still prefer a safe points lead to a risky attack.
\end{itemize}

\subsection{Verdict and next step}
\textbf{Adopted} as the metric of all later experiments. The champion is unchanged. Next: why do games stall? (experiment 6).

\subsection{Notes for the paper}
Draft: evaluation-metric section. Points share at a turn cap is blind to stalling; the finish score separates settings that points share ranks as equal. Candidate figure: finish, points, conquest rate and turns against \knob{front} from 0.5 to 32.

\todo{jakob}{Is this worth a place in the paper? Where?}


\clearpage\phantomsection\addcontentsline{toc}{section}{Experiment 6: Endgame peace break (kept)}

% generated by scripts/experiment-report.ts from experiment.json and models/<run>/ - do not edit
\providecommand{\ExpHeader}{}\renewcommand{\ExpHeader}{%
{\sffamily\footnotesize\color{expgrey}EXPERIMENT 6 \textperiodcentered\ MECHANISM \hfill 2026-11-24}\par
{\LARGE\bfseries\sffamily Endgame peace break}\hfill\verdictbox{expgreen}{KEPT}\par\smallskip
{\itshape Do games stall because the last players are at peace, and does breaking that peace make the bot finish?}\par\medskip
\begin{tabular}{@{}p{22mm}p{\dimexpr\linewidth-26mm}@{}}
\fact{Change}{Break peace with a partner when every living rival is a peace partner and we hold at least breakPeace times its cities.}
\fact{Knobs}{\knob{breakPeace} $\in$ \{0.5, 1, 2, 4\}, chosen 0.5}
\fact{Base}{\model{greedy-full-front} $\rightarrow$ result \model{greedy-full-peacebreak}}
\fact{Method}{Diagnosis of stalled games, then strength sweeps in two setups (200 games per point, final test 400); metric: finish score}
\fact{Setups}{4 players, 16$\times$16, drylands, turn cap 150; 8 players, 20$\times$20, drylands, turn cap 300}
\fact{Record}{issue \issue{124}, pull request \issue{129}; status: done}
\end{tabular}\par\medskip
\begin{center}\begin{tikzpicture}[node distance=5mm]
    \node[flowbox] (n0) {{\sffamily\bfseries\scriptsize\color{expgrey}BASE}\\[1pt]\model{full-front}};
    \node[flowbox,right=5mm of n0] (n1) {{\sffamily\bfseries\scriptsize\color{expgrey}CHANGE}\\[1pt]\knob{breakPeace}\\{\scriptsize 0.5, 1, 2, 4}};
    \node[flowbox,right=5mm of n1] (n2) {{\sffamily\bfseries\scriptsize\color{expgrey}MEASURED}\\[1pt]Diagnosis, then sweeps in 2 setups\\{\scriptsize 2 setups, metric: finish score}};
    \node[flowbox,right=5mm of n2] (n3) {{\sffamily\bfseries\scriptsize\color{expgrey}RESULT}\\[1pt]Finish score, 6-model panel, 4 players\\{\scriptsize 75.0 $\rightarrow$ 78.8\,\%}\\\textbf{+3.8\pp}};
    \node[flowbox,right=5mm of n3] (n4) {{\sffamily\bfseries\scriptsize\color{expgrey}VERDICT}\\[1pt]\textbf{\color{expgreen}KEPT}\\{\scriptsize new champion:}\\\model{full-peacebreak}};
    \draw[flowarrow] (n0) -- (n1); \draw[flowarrow] (n1) -- (n2); \draw[flowarrow] (n2) -- (n3); \draw[flowarrow] (n3) -- (n4);
\end{tikzpicture}\end{center}
\begin{center}\sffamily\small\begin{tabular}{@{}lrrr@{}}\toprule
Measure & Base & New & Change\\\midrule
Finish score, 6-model panel, 4 players & 75.0\,\% & 78.8\,\% & +3.8\\
Finish score, 6-model panel, 8 players & 77.4\,\% & 81.0\,\% & +3.6\\
Conquest wins, 4 players & 64.0\,\% & 76.5\,\% & +12.5\\
Game length in turns, 4 players & 75.3 & 62.9 & $-$12.4\\
Points vs base, final test & 50.0\,\% & 50.4\,\% & +0.4\\
\bottomrule\end{tabular}\end{center}
}
\providecommand{\ExpResults}{}\renewcommand{\ExpResults}{%
\subsection{Sweep, 4 players {\normalfont\small(\model{greedy-full-peacebreak})}}
\begin{center}\sffamily\small\begin{tabular}{@{}lrrrrrr@{}}\toprule
Point & Finish (\%) & Points (\%) & Won (\%) & Conquest (\%) & Turns & vs base: points (\%)\\\midrule
base & 72.43 & 67.36 & 80.7 & 64.4 & 69.8 & 50.0\\
breakPeace=0.5 & 74.64 & 67.93 & 80.6 & 72.5 & 61.1 & 50.0\\
breakPeace=1 & 74.57 & 67.84 & 80.2 & 72.1 & 61.7 & 50.3\\
breakPeace=2 & 74.13 & 67.64 & 80.7 & 70.1 & 64.2 & 50.2\\
breakPeace=4 & 73.27 & 67.40 & 80.7 & 67.2 & 67.0 & 49.9\\
\bottomrule\end{tabular}\\[3pt]{\footnotesize\color{expgrey}Panel means; 200 games against each of 7 panel models per point.}\end{center}
\begin{center}\begin{tikzpicture}\begin{axis}[expplot,height=4.6cm,xtick={0,1,2,3,4},xticklabels={{base},{0.5},{1},{2},{4}},
  xlabel={\texttt{breakPeace}},ylabel={Finish score, panel mean (\%)},enlarge x limits=0.12,enlarge y limits=0.25]
\addplot[dashed,expgrey,line width=0.7pt,forget plot] coordinates {(0,72.429) (4,72.429)};
\addplot[only marks,mark=*,mark size=2.6pt,expblue] coordinates { (1,74.643) (2,74.571) (3,74.129) (4,73.271)};
\addplot[only marks,mark=*,mark size=2.6pt,expgrey] coordinates {(0,72.429)};
\end{axis}\end{tikzpicture}\\[-2pt]{\footnotesize\color{expgrey}Dashed: the base. Grey: base; blue: sweep points.}\end{center}
\begin{center}\sffamily\small\begin{tabular}{@{}lrrrrr@{}}\toprule
Opponent & Won (\%) & Points (\%) & Finish (\%) & Conquest (\%) & Turns\\\midrule
\model{v1} & 96 $\pm$ 2 & 77.9 & 87.7 & 90 & 54\\
\model{hill-v1} & 91 $\pm$ 3 & 73.6 & 82.2 & 80 & 63\\
\model{hill-v2} & 95 $\pm$ 2 & 75.0 & 84.8 & 83 & 65\\
\model{hill-v3} & 91 $\pm$ 3 & 72.3 & 83.4 & 84 & 61\\
\model{hill-v4} & 88 $\pm$ 3 & 69.8 & 77.9 & 71 & 72\\
\model{full-v1} & 57 $\pm$ 5 & 54.7 & 56.7 & 52 & 63\\
\model{full-front} & 51 $\pm$ 5 & 50.4 & 51.7 & 47 & 61\\
\midrule
\textbf{Mean without \model{full-front}} & 86.2 & 70.5 & 78.8 & 76.5 & 62.9\\
\bottomrule\end{tabular}\\[3pt]{\footnotesize\color{expgrey}Final test of breakPeace=0.5: 400 unseen games against each.}\end{center}

\subsection{Sweep, 8 players {\normalfont\small(\model{greedy-full-peacebreak-big})}}
\begin{center}\sffamily\small\begin{tabular}{@{}lrrrrrr@{}}\toprule
Point & Finish (\%) & Points (\%) & Won (\%) & Conquest (\%) & Turns & vs base: points (\%)\\\midrule
base & 74.04 & 71.31 & 83.3 & 57.9 & 88.1 & 50.0\\
breakPeace=0.5 & 78.00 & 73.21 & 84.1 & 69.6 & 78.3 & 50.4\\
breakPeace=1 & 77.47 & 72.76 & 83.4 & 68.1 & 79.7 & 50.4\\
breakPeace=2 & 76.31 & 72.01 & 83.3 & 64.9 & 82.2 & 50.4\\
breakPeace=4 & 75.01 & 71.47 & 83.3 & 61.1 & 85.4 & 50.0\\
\bottomrule\end{tabular}\\[3pt]{\footnotesize\color{expgrey}Panel means; 200 games against each of 7 panel models per point.}\end{center}
\begin{center}\begin{tikzpicture}\begin{axis}[expplot,height=4.6cm,xtick={0,1,2,3,4},xticklabels={{base},{0.5},{1},{2},{4}},
  xlabel={\texttt{breakPeace}},ylabel={Finish score, panel mean (\%)},enlarge x limits=0.12,enlarge y limits=0.25]
\addplot[dashed,expgrey,line width=0.7pt,forget plot] coordinates {(0,74.043) (4,74.043)};
\addplot[only marks,mark=*,mark size=2.6pt,expblue] coordinates { (1,78.000) (2,77.471) (3,76.314) (4,75.014)};
\addplot[only marks,mark=*,mark size=2.6pt,expgrey] coordinates {(0,74.043)};
\end{axis}\end{tikzpicture}\\[-2pt]{\footnotesize\color{expgrey}Dashed: the base. Grey: base; blue: sweep points.}\end{center}
\begin{center}\sffamily\small\begin{tabular}{@{}lrrrrr@{}}\toprule
Opponent & Won (\%) & Points (\%) & Finish (\%) & Conquest (\%) & Turns\\\midrule
\model{v1} & 97 $\pm$ 2 & 84.0 & 89.5 & 84 & 70\\
\model{hill-v1} & 96 $\pm$ 2 & 81.8 & 85.9 & 74 & 81\\
\model{hill-v2} & 97 $\pm$ 2 & 82.0 & 85.8 & 73 & 83\\
\model{hill-v3} & 92 $\pm$ 3 & 76.6 & 84.2 & 75 & 81\\
\model{hill-v4} & 95 $\pm$ 2 & 77.7 & 82.1 & 66 & 89\\
\model{full-v1} & 60 $\pm$ 5 & 56.1 & 58.7 & 47 & 85\\
\model{full-front} & 50 $\pm$ 5 & 50.1 & 50.7 & 37 & 88\\
\midrule
\textbf{Mean without \model{full-front}} & 89.3 & 76.4 & 81.0 & 69.7 & 81.5\\
\bottomrule\end{tabular}\\[3pt]{\footnotesize\color{expgrey}Final test of breakPeace=0.5: 400 unseen games against each.}\end{center}

\subsection{The base on the final-test maps, 4 players {\normalfont\small(\model{greedy-full-front-final})}}
\begin{center}\sffamily\small\begin{tabular}{@{}lrrrrr@{}}\toprule
Opponent & Won (\%) & Points (\%) & Finish (\%) & Conquest (\%) & Turns\\\midrule
\model{v1} & 94 $\pm$ 2 & 75.7 & 83.1 & 76 & 68\\
\model{hill-v1} & 91 $\pm$ 3 & 72.3 & 78.3 & 68 & 75\\
\model{hill-v2} & 95 $\pm$ 2 & 73.7 & 80.9 & 71 & 77\\
\model{hill-v3} & 91 $\pm$ 3 & 70.3 & 78.3 & 68 & 76\\
\model{hill-v4} & 87 $\pm$ 3 & 68.6 & 74.1 & 59 & 83\\
\model{full-v1} & 58 $\pm$ 5 & 54.8 & 55.1 & 44 & 73\\
\model{full-front} & 50 $\pm$ 5 & 50.0 & 50.0 & 40 & 69\\
\midrule
\textbf{Mean without \model{full-front}} & 85.9 & 69.2 & 75.0 & 64.0 & 75.3\\
\bottomrule\end{tabular}\\[3pt]{\footnotesize\color{expgrey}Final test of greedy-full-front: 400 unseen games against each.}\end{center}

\subsection{The base on the final-test maps, 8 players {\normalfont\small(\model{greedy-full-front-final-big})}}
\begin{center}\sffamily\small\begin{tabular}{@{}lrrrrr@{}}\toprule
Opponent & Won (\%) & Points (\%) & Finish (\%) & Conquest (\%) & Turns\\\midrule
\model{v1} & 96 $\pm$ 2 & 81.9 & 85.5 & 72 & 80\\
\model{hill-v1} & 96 $\pm$ 2 & 80.8 & 83.8 & 68 & 86\\
\model{hill-v2} & 96 $\pm$ 2 & 80.0 & 81.3 & 60 & 94\\
\model{hill-v3} & 90 $\pm$ 3 & 73.5 & 78.0 & 58 & 94\\
\model{hill-v4} & 95 $\pm$ 2 & 75.8 & 78.3 & 56 & 98\\
\model{full-v1} & 59 $\pm$ 5 & 55.5 & 57.3 & 42 & 89\\
\model{full-front} & 50 $\pm$ 5 & 50.0 & 50.0 & 32 & 93\\
\midrule
\textbf{Mean without \model{full-front}} & 88.7 & 74.6 & 77.4 & 59.3 & 90.2\\
\bottomrule\end{tabular}\\[3pt]{\footnotesize\color{expgrey}Final test of greedy-full-front: 400 unseen games against each.}\end{center}
}

\ExpHeader

\subsection{Question and hypothesis}
Experiment 5 showed that 19\,\% (4 players) to 36\,\% (8 players) of the champion's self-play games never finish. Opening four stalled games at turn 300: in two, the last players were at peace; in the other two a giant army sat idle (experiment 7). The greedy bot scores a peace break at $-1$, so it never breaks peace, and a game between peace partners cannot end. \textbf{Hypothesis:} breaking peace once nobody else is left to fight ends these games and raises the finish score.

\subsection{Mechanism}
The action \texttt{BreakPeace} with partner $p$ scores 100 when both hold: (1) every living rival is a peace partner, and (2) $\mathit{cities}(\text{me}) \ge \knob{breakPeace}\cdot\mathit{cities}(p)$. Otherwise it keeps its score of $-1$. Default 0 (off). Code: \texttt{endgame\_break} in \texttt{bot\_full.rs}.

\subsection{Setup}
Sanity check in self-play (40 maps): finished games rise from 31 to 34 (4 players) and from 21 to 30 (8 players). The sweep ran in both setups. The base was measured again on the final-test maps (the two extra runs below), so that the 6-model panel means of base and new model come from the same maps.

\subsection{Results}
\ExpResults

The best value is 0.5 in both setups, and the effect shrinks as the ratio rises (4 is about the base). On the final-test maps against the six fixed panel models, the finish score rises by 3.8 points (4 players) and 3.6 (8 players); conquest wins go from 64.0 to 76.5\,\% and games are 12 turns shorter. Chance moves these 6-model means by about $\pm$1.2. Head to head against the base it is a draw: finish 51.7\,\%, points 50.4\,\%, won 51\,\% ($\pm$5).

\subsection{Behaviour}
Self-play games finished by the turn cap: 67\,\% to 78\,\% (8 players), 82\,\% to 94\,\% (4 players).

\subsection{Interpretation}
The break only fires when the last rivals are at peace. In a match between two models that almost never happens between the two of them, so the head-to-head cannot see the mechanism; against a panel of varied opponents it shows clearly. A value of 0.5 means the bot attacks a partner holding up to twice its cities: ending the game matters more than waiting for an advantage.

\todo{jakob}{Your reading: does this match what you saw? Anything to add or correct?}

\subsection{Threats to validity}
\begin{itemize}
\item The keep rule was changed for this experiment (the panel rule), after seeing the result. It is principled (the head-to-head is blind here) but it was not pre-registered.
\item Opponents never break peace themselves; against bots or humans that do, the gain may be smaller.
\item The finish score rewards finishing by construction, so part of the gain is the metric agreeing with the mechanism. Points share rose too (+1.3 and +1.8), which is the independent evidence.
\end{itemize}

\subsection{Verdict and next step}
\textbf{Kept} under the panel keep rule (finish on the 6-model panel more than 1.5 points above the base, head-to-head not worse), added with Jakob's go on 24 November. New champion: \model{greedy-full-peacebreak}. Next: the idle giant armies (experiment 7).

\subsection{Notes for the paper}
Draft: (1) a diplomacy rule that never breaks peace makes a fifth to a third of self-play games unwinnable; one rule fixes most of it. (2) Methodological: pairwise head-to-head tests are blind to mechanisms that act on third parties; a fixed opponent panel is needed. Candidate table: the final-test comparison in both setups.

\todo{jakob}{Is this worth a place in the paper? Where?}


\clearpage\phantomsection\addcontentsline{toc}{section}{Experiment 7: Movement noise (kept)}

% generated by scripts/experiment-report.ts from experiment.json and models/<run>/ - do not edit
\providecommand{\ExpHeader}{}\renewcommand{\ExpHeader}{%
{\sffamily\footnotesize\color{expgrey}EXPERIMENT 7 \textperiodcentered\ MECHANISM \hfill 2026-11-25}\par
{\LARGE\bfseries\sffamily Movement noise}\hfill\verdictbox{expgreen}{KEPT}\par\smallskip
{\itshape Do jammed armies come loose when moves get a zero-mean random term and stuck units may step sideways?}\par\medskip
\begin{tabular}{@{}p{22mm}p{\dimexpr\linewidth-26mm}@{}}
\fact{Change}{Every move gets a random term uniform in [$-$moveNoise, +moveNoise], with mean zero; a move that gains nothing toward the goal is no longer ruled out and is taken when its score is positive.}
\fact{Knobs}{\knob{moveNoise} $\in$ \{1, 2, 4, 8\}, chosen 4}
\fact{Base}{\model{greedy-full-peacebreak} $\rightarrow$ result \model{greedy-full-movenoise}}
\fact{Method}{Diagnosis of stalled games, then strength sweeps in two setups (200 games per point, final test 400); metric: finish score}
\fact{Setups}{4 players, 16$\times$16, drylands, turn cap 150; 8 players, 20$\times$20, drylands, turn cap 300}
\fact{Record}{issue \issue{130}; status: running}
\end{tabular}\par\medskip
\begin{center}\begin{tikzpicture}[node distance=5mm]
    \node[flowbox] (n0) {{\sffamily\bfseries\scriptsize\color{expgrey}BASE}\\[1pt]\model{full-peacebreak}};
    \node[flowbox,right=5mm of n0] (n1) {{\sffamily\bfseries\scriptsize\color{expgrey}CHANGE}\\[1pt]\knob{moveNoise}\\{\scriptsize 1, 2, 4, 8}};
    \node[flowbox,right=5mm of n1] (n2) {{\sffamily\bfseries\scriptsize\color{expgrey}MEASURED}\\[1pt]Diagnosis, then sweeps in 2 setups\\{\scriptsize 2 setups, metric: finish score}};
    \node[flowbox,right=5mm of n2] (n3) {{\sffamily\bfseries\scriptsize\color{expgrey}RESULT}\\[1pt]Finish vs base, final test\\{\scriptsize 50.0 $\rightarrow$ 60.1\,\%}\\\textbf{+10.1\pp}};
    \node[flowbox,right=5mm of n3] (n4) {{\sffamily\bfseries\scriptsize\color{expgrey}VERDICT}\\[1pt]\textbf{\color{expgreen}KEPT}\\{\scriptsize new champion:}\\\model{full-movenoise}};
    \draw[flowarrow] (n0) -- (n1); \draw[flowarrow] (n1) -- (n2); \draw[flowarrow] (n2) -- (n3); \draw[flowarrow] (n3) -- (n4);
\end{tikzpicture}\end{center}
\begin{center}\sffamily\small\begin{tabular}{@{}lrrr@{}}\toprule
Measure & Base & New & Change\\\midrule
Finish vs base, final test & 50.0\,\% & 60.1\,\% & +10.1\\
Won vs base, final test & 50.0\,\% & 62.5\,\% & +12.5\\
Points vs base, final test & 50.0\,\% & 56.9\,\% & +6.9\\
Finish score, 6-model panel, 4 players & 78.8\,\% & 81.3\,\% & +2.5\\
Conquest wins, 4 players & 76.5\,\% & 80.5\,\% & +4.0\\
Game length in turns, 4 players & 62.9 & 58.6 & $-$4.3\\
\bottomrule\end{tabular}\end{center}
}
\providecommand{\ExpResults}{}\renewcommand{\ExpResults}{%
\subsection{Sweep, 4 players {\normalfont\small(\model{greedy-full-movenoise})}}
\begin{center}\sffamily\small\begin{tabular}{@{}lrrrrrr@{}}\toprule
Point & Finish (\%) & Points (\%) & Won (\%) & Conquest (\%) & Turns & vs base: points (\%)\\\midrule
base & 74.46 & 67.93 & 80.6 & 72.5 & 60.5 & 50.0\\
moveNoise=1 & 78.37 & 69.84 & 85.1 & 76.4 & 58.1 & 52.3\\
moveNoise=2 & 77.64 & 70.10 & 84.9 & 75.2 & 58.6 & 52.7\\
moveNoise=4 & 79.21 & 71.40 & 86.3 & 78.9 & 56.4 & 55.7\\
moveNoise=8 & 78.11 & 70.90 & 84.4 & 78.5 & 56.6 & 53.0\\
\bottomrule\end{tabular}\\[3pt]{\footnotesize\color{expgrey}Panel means; 200 games against each of 7 panel models per point.}\end{center}
\begin{center}\begin{tikzpicture}\begin{axis}[expplot,height=4.6cm,xtick={0,1,2,3,4},xticklabels={{base},{1},{2},{4},{8}},
  xlabel={\texttt{moveNoise}},ylabel={Finish score, panel mean (\%)},enlarge x limits=0.12,enlarge y limits=0.25]
\addplot[dashed,expgrey,line width=0.7pt,forget plot] coordinates {(0,74.457) (4,74.457)};
\addplot[only marks,mark=*,mark size=2.6pt,expblue] coordinates { (1,78.371) (2,77.643) (3,79.214) (4,78.114)};
\addplot[only marks,mark=*,mark size=2.6pt,expgrey] coordinates {(0,74.457)};
\end{axis}\end{tikzpicture}\\[-2pt]{\footnotesize\color{expgrey}Dashed: the base. Grey: base; blue: sweep points.}\end{center}
\begin{center}\sffamily\small\begin{tabular}{@{}lrrrrr@{}}\toprule
Opponent & Won (\%) & Points (\%) & Finish (\%) & Conquest (\%) & Turns\\\midrule
\model{v1} & 97 $\pm$ 2 & 79.5 & 88.8 & 92 & 49\\
\model{hill-v1} & 91 $\pm$ 3 & 75.3 & 84.0 & 84 & 58\\
\model{hill-v2} & 93 $\pm$ 3 & 75.0 & 85.2 & 85 & 59\\
\model{hill-v3} & 90 $\pm$ 3 & 72.3 & 82.7 & 82 & 60\\
\model{hill-v4} & 94 $\pm$ 2 & 73.8 & 82.9 & 78 & 65\\
\model{full-v1} & 67 $\pm$ 5 & 59.9 & 64.1 & 62 & 60\\
\model{full-peacebreak} & 63 $\pm$ 5 & 56.9 & 60.1 & 56 & 59\\
\midrule
\textbf{Mean without \model{full-peacebreak}} & 88.6 & 72.6 & 81.3 & 80.5 & 58.6\\
\bottomrule\end{tabular}\\[3pt]{\footnotesize\color{expgrey}Final test of moveNoise=4: 400 unseen games against each.}\end{center}

\subsection{Sweep, 8 players {\normalfont\small(\model{greedy-full-movenoise-big})}}
{\small\color{expgrey}No points measured yet.}
}

\ExpHeader

\subsection{Question and hypothesis}
Jakob, 24 November: ``what about adding stochastic movement with variance parameter but fixed expectation value the movement weights from the champion to get around obstacles and also units spreading out to a degree which can be hypertuned on top champion like density parameter''.

Diagnosis first (six stalled 8-player self-play games of the champion, opened at turn 250): the leader had 97 to 142 units, of which 53 to 77 could not move at all and most of the rest had no move toward a goal; only 0 to 2 units per game had a useful move. Between 60 and 105 of them had six or more own units as neighbours. The armies are jammed by their own units. \textbf{Hypothesis:} a random term on move scores with mean zero keeps the champion's preferences on average but lets a stuck unit step aside and find another route, so more games finish.

The prediction was not written down as numbers before the first sweep (this notebook did not exist yet). The self-play check before the sweep gave: finished games 34 to 37 or 39 of 40 with 4 players, but 30 to 25 or 29 of 40 with 8 players.

\subsection{Mechanism}
Let $g$ be a move's gain toward the unit's goal in tiles and $z \sim U(-\knob{moveNoise}, +\knob{moveNoise})$, seeded by game, turn, unit and target tile, so a game stays reproducible. For $g>0$ the move scores as before plus $z$. For $g\le 0$ (sideways or back), which the base rules out with $-1$, the move scores $g + z$ and is taken when that is positive. The bot already had a small noise term (\knob{greedy.noise} $=1$), but it has mean $+0.5$ and applies only to moves with $g>0$, so it cannot unjam anything. Default 0. Code: \texttt{move\_score} in \texttt{bot\_full.rs}.

\subsection{Setup}
Values sized against a move score of about 25 plus the tiles gained: 1, 2, 4, 8. The seven-model panel is the fixed six plus the base.

\subsection{Results}
\ExpResults

\textbf{4 players.} Every tested value beats the base in the sweep, on the panel and head to head; 4 is best (79.2\,\% finish on seven models against 74.5\,\%). In the final test \knob{moveNoise} $=4$ beats the champion clearly: finish 60.1\,\%, points 56.9\,\%, won 62\,\% ($\pm$5). Against the fixed six on the final-test maps the finish score rises from 78.8 to 81.3\,\% (the base's figure is from experiment 6, same maps), conquest wins from 76.5 to 80.5\,\%, and games are four turns shorter.

\textbf{8 players.} The same picture on 20$\times$20 maps: every value beats the base, 4 and 8 most. In the final test \knob{moveNoise} $=4$ scores finish 59.4\,\%, points 57.6\,\%, won 61\,\% ($\pm$5) against the champion, and 83.7\,\% against 81.0\,\% on the fixed six (the base's figure is from experiment 6, same maps), with conquest wins up from 69.7 to 73.4\,\% and games four turns shorter. The self-play check before the sweep had looked mixed in this setup (25 to 29 of 40 games finished, against 30); the panel measurement does not bear that out.

\subsection{Behaviour}
The six stalled 8-player games of the diagnosis, replayed with \knob{moveNoise} $=4$: four now finish, at turns 50 to 75; none did before. In the two that still stall, the leader's army is as jammed as before (55\,\% of its units cannot move or attack, against 42\,\% for the champion over all six; 72\,\% have six or more own neighbours). The noise keeps most jams from forming; it does not dissolve one that has formed.

\todo{jakob}{What you saw in the live games, if you watched any.}

\subsection{Interpretation}
The gain is far larger than the late-game unjamming alone would give: the replayed games end at turns 50 to 75, not at turn 280. Allowing sideways steps changes play from the start, because a unit behind another unit is no longer frozen until the one in front moves. The response is flat from 1 to 8, as with the front pull (experiment 3): what matters is that the step is allowed at all, and any noise large enough to make $g+z$ positive now and then does that. It is the largest head-to-head gain of the series so far.

One caveat on the design: the knob does two things at once, zero-mean noise on all moves and permission for moves with $g\le 0$. The flat response suggests the permission carries the effect, but the two were not separated.

\todo{jakob}{Your reading: agree, disagree, what is missing?}

\subsection{Threats to validity}
\begin{itemize}
\item Noise and the sideways permission are confounded in one knob (see above). A variant with permission but no noise would separate them.
\item The opponents are all greedy bots that jam the same way; part of the gain is exploiting their jams.
\item The prediction was not fixed as numbers before the first sweep.
\item Self-play finish counts on 40 maps (the check before the sweep) pointed the wrong way for 8 players; such small checks are useful for sizing a knob, not for judging it.
\end{itemize}

\subsection{Verdict and next step}
\textbf{Kept} under the head-to-head rule (finish and points above 50\,\%, the $\pm$5 interval clear of 50\,\%, panel mean above the base's). New champion: \model{greedy-full-movenoise}. Next: spread on top of it (experiment 8); then formations (\issue{132}) and the same knobs for the plain greedy bots (\issue{131}).

\subsection{Notes for the paper}
Draft: a deterministic greedy policy deadlocks on its own units; a small zero-mean perturbation with permission for non-improving moves removes most deadlocks and is worth about 60\,\% head to head. Candidate figure: finished games and game length for the six diagnosed seeds, base against new; plus the flat sweep response.

\todo{jakob}{Is this worth a place in the paper? Where?}


\clearpage\phantomsection\addcontentsline{toc}{section}{Experiment 8: Spread out: thin own crowds (not kept)}

% generated by scripts/experiment-report.ts from experiment.json and models/<run>/ - do not edit
\providecommand{\ExpHeader}{}\renewcommand{\ExpHeader}{%
{\sffamily\footnotesize\color{expgrey}EXPERIMENT 8 \textperiodcentered\ MECHANISM \hfill 2026-11-25}\par
{\LARGE\bfseries\sffamily Spread out: thin own crowds}\hfill\verdictbox{expred}{NOT KEPT}\par\smallskip
{\itshape Do armies jam less when a move is rewarded for leaving a crowd of own units?}\par\medskip
\begin{tabular}{@{}p{22mm}p{\dimexpr\linewidth-26mm}@{}}
\fact{Change}{A move gets +spread per own unit fewer next to the target tile than next to the unit's current tile; like movement noise, it also allows a sideways step when that makes the score positive.}
\fact{Knobs}{\knob{spread} $\in$ \{0.25, 0.5, 1, 2, 4\}, best 0.5}
\fact{Base}{\model{greedy-full-peacebreak}}
\fact{Method}{Strength sweeps (200 games per point, final test 400): on the champion, and on the champion plus movement noise; metric: finish score}
\fact{Setups}{4 players, 16$\times$16, drylands, turn cap 150}
\fact{Record}{issue \issue{130}; status: done}
\end{tabular}\par\medskip
\begin{center}\begin{tikzpicture}[node distance=5mm]
    \node[flowbox] (n0) {{\sffamily\bfseries\scriptsize\color{expgrey}BASE}\\[1pt]\model{full-peacebreak}};
    \node[flowbox,right=5mm of n0] (n1) {{\sffamily\bfseries\scriptsize\color{expgrey}CHANGE}\\[1pt]\knob{spread}\\{\scriptsize 0.25, 0.5, 1, 2, 4}};
    \node[flowbox,right=5mm of n1] (n2) {{\sffamily\bfseries\scriptsize\color{expgrey}MEASURED}\\[1pt]Sweeps on two bases\\{\scriptsize 1 setup, metric: finish score}};
    \node[flowbox,right=5mm of n2] (n3) {{\sffamily\bfseries\scriptsize\color{expgrey}RESULT}\\[1pt]Finish score, 6-model panel, on top of movement noise\\{\scriptsize 81.3 $\rightarrow$ 81.6\,\%}\\\textbf{+0.3\pp}};
    \node[flowbox,right=5mm of n3] (n4) {{\sffamily\bfseries\scriptsize\color{expgrey}VERDICT}\\[1pt]\textbf{\color{expred}NOT KEPT}};
    \draw[flowarrow] (n0) -- (n1); \draw[flowarrow] (n1) -- (n2); \draw[flowarrow] (n2) -- (n3); \draw[flowarrow] (n3) -- (n4);
\end{tikzpicture}\end{center}
\begin{center}\sffamily\small\begin{tabular}{@{}lrrr@{}}\toprule
Measure & Base & New & Change\\\midrule
Finish score, 6-model panel, on top of movement noise & 81.3\,\% & 81.6\,\% & +0.3\\
Finish vs base, final test, on top of movement noise & 50.0\,\% & 48.8\,\% & $-$1.2\\
Finish score, 6-model panel, alone on the champion & 78.8\,\% & 79.9\,\% & +1.1\\
Finish vs base, final test, alone on the champion & 50.0\,\% & 54.7\,\% & +4.7\\
\bottomrule\end{tabular}\end{center}
}
\providecommand{\ExpResults}{}\renewcommand{\ExpResults}{%
\subsection{Sweep on the champion (peacebreak) {\normalfont\small(\model{greedy-full-spread})}}
\begin{center}\sffamily\small\begin{tabular}{@{}lrrrrrr@{}}\toprule
Point & Finish (\%) & Points (\%) & Won (\%) & Conquest (\%) & Turns & vs base: points (\%)\\\midrule
base & 74.46 & 67.93 & 80.6 & 72.5 & 60.5 & 50.0\\
spread=0.5 & 79.10 & 70.09 & 85.9 & 77.3 & 56.2 & 54.3\\
spread=1 & 78.23 & 69.46 & 84.4 & 76.1 & 55.8 & 54.3\\
spread=2 & 77.60 & 69.34 & 83.8 & 76.1 & 54.7 & 55.4\\
spread=4 & 77.54 & 68.46 & 82.2 & 77.8 & 54.2 & 52.1\\
\bottomrule\end{tabular}\\[3pt]{\footnotesize\color{expgrey}Panel means; 200 games against each of 7 panel models per point.}\end{center}
\begin{center}\begin{tikzpicture}\begin{axis}[expplot,height=4.6cm,xtick={0,1,2,3,4},xticklabels={{base},{0.5},{1},{2},{4}},
  xlabel={\texttt{spread}},ylabel={Finish score, panel mean (\%)},enlarge x limits=0.12,enlarge y limits=0.25]
\addplot[dashed,expgrey,line width=0.7pt,forget plot] coordinates {(0,74.457) (4,74.457)};
\addplot[only marks,mark=*,mark size=2.6pt,expblue] coordinates { (1,79.100) (2,78.229) (3,77.600) (4,77.543)};
\addplot[only marks,mark=*,mark size=2.6pt,expgrey] coordinates {(0,74.457)};
\end{axis}\end{tikzpicture}\\[-2pt]{\footnotesize\color{expgrey}Dashed: the base. Grey: base; blue: sweep points.}\end{center}
\begin{center}\sffamily\small\begin{tabular}{@{}lrrrrr@{}}\toprule
Opponent & Won (\%) & Points (\%) & Finish (\%) & Conquest (\%) & Turns\\\midrule
\model{v1} & 95 $\pm$ 2 & 77.1 & 87.0 & 88 & 50\\
\model{hill-v1} & 91 $\pm$ 3 & 73.3 & 83.0 & 81 & 56\\
\model{hill-v2} & 93 $\pm$ 3 & 74.0 & 84.7 & 83 & 58\\
\model{hill-v3} & 91 $\pm$ 3 & 71.3 & 83.4 & 81 & 61\\
\model{hill-v4} & 91 $\pm$ 3 & 70.7 & 80.5 & 74 & 66\\
\model{full-v1} & 63 $\pm$ 5 & 56.4 & 60.6 & 57 & 59\\
\model{full-peacebreak} & 55 $\pm$ 5 & 52.2 & 54.7 & 52 & 56\\
\midrule
\textbf{Mean without \model{full-peacebreak}} & 87.1 & 70.5 & 79.9 & 77.2 & 58.5\\
\bottomrule\end{tabular}\\[3pt]{\footnotesize\color{expgrey}Final test of spread=0.5: 400 unseen games against each.}\end{center}

\subsection{Sweep on top of movement noise 4 {\normalfont\small(\model{greedy-full-noise-spread})}}
\begin{center}\sffamily\small\begin{tabular}{@{}lrrrrrr@{}}\toprule
Point & Finish (\%) & Points (\%) & Won (\%) & Conquest (\%) & Turns & vs base: points (\%)\\\midrule
base & 78.11 & 70.59 & 84.9 & 77.5 & 56.5 & 50.0\\
spread=0.25 & 77.93 & 70.06 & 83.9 & 76.8 & 56.5 & 51.0\\
spread=0.5 & 77.97 & 70.03 & 84.4 & 77.1 & 56.0 & 50.3\\
spread=1 & 77.80 & 68.96 & 83.8 & 77.6 & 55.1 & 48.3\\
spread=2 & 77.49 & 69.30 & 83.3 & 77.6 & 54.2 & 49.1\\
\bottomrule\end{tabular}\\[3pt]{\footnotesize\color{expgrey}Panel means; 200 games against each of 7 panel models per point.}\end{center}
\begin{center}\begin{tikzpicture}\begin{axis}[expplot,height=4.6cm,xtick={0,1,2,3,4},xticklabels={{base},{0.25},{0.5},{1},{2}},
  xlabel={\texttt{spread}},ylabel={Finish score, panel mean (\%)},enlarge x limits=0.12,enlarge y limits=0.25]
\addplot[dashed,expgrey,line width=0.7pt,forget plot] coordinates {(0,78.114) (4,78.114)};
\addplot[only marks,mark=*,mark size=2.6pt,expblue] coordinates { (1,77.929) (2,77.971) (3,77.800) (4,77.486)};
\addplot[only marks,mark=*,mark size=2.6pt,expgrey] coordinates {(0,78.114)};
\end{axis}\end{tikzpicture}\\[-2pt]{\footnotesize\color{expgrey}Dashed: the base. Grey: base; blue: sweep points.}\end{center}
\begin{center}\sffamily\small\begin{tabular}{@{}lrrrrr@{}}\toprule
Opponent & Won (\%) & Points (\%) & Finish (\%) & Conquest (\%) & Turns\\\midrule
\model{v1} & 96 $\pm$ 2 & 79.7 & 88.4 & 91 & 50\\
\model{hill-v1} & 92 $\pm$ 3 & 76.0 & 84.9 & 87 & 54\\
\model{hill-v2} & 94 $\pm$ 2 & 75.6 & 85.7 & 87 & 57\\
\model{hill-v3} & 92 $\pm$ 3 & 72.6 & 84.8 & 87 & 53\\
\model{hill-v4} & 91 $\pm$ 3 & 71.5 & 81.5 & 77 & 65\\
\model{full-v1} & 68 $\pm$ 5 & 60.3 & 64.3 & 63 & 58\\
\model{full-movenoise} & 49 $\pm$ 5 & 49.2 & 48.8 & 44 & 57\\
\midrule
\textbf{Mean without \model{full-movenoise}} & 88.6 & 72.6 & 81.6 & 81.7 & 56.0\\
\bottomrule\end{tabular}\\[3pt]{\footnotesize\color{expgrey}Final test of spread=0.5: 400 unseen games against each.}\end{center}
}

\ExpHeader

\subsection{Question and hypothesis}
Jakob, 24 November: ``what about adding stochastic movement with variance parameter but fixed expectation value the movement weights from the champion to get around obstacles and also units spreading out to a degree which can be hypertuned on top champion like density parameter''.

Diagnosis first (six stalled 8-player self-play games of the champion, opened at turn 250): the leader had 97 to 142 units, of which 53 to 77 could not move at all and most of the rest had no move toward a goal; only 0 to 2 units per game had a useful move. Between 60 and 105 of them had six or more own units as neighbours. The armies are jammed by their own units. \textbf{Hypothesis:} a bonus for moving to a tile with fewer own neighbours thins the crowds, so fewer units are boxed in and more games finish. Open question: does it add anything once movement noise (experiment 7) is on, since both act on the same jam?

\subsection{Mechanism}
Let $c(x)$ be the number of own units within one tile of $x$. A move from $a$ to $b$ gets $+\,\knob{spread}\cdot(c(a)-c(b))$. As in experiment 7, a move with no gain toward the goal is taken when this makes its score positive. Default 0. Code: \texttt{move\_score} in \texttt{bot\_full.rs}.

\subsection{Setup}
One mechanism per experiment: the knob is first swept alone on the champion of experiment 6, then on top of \knob{moveNoise} $=4$ to see whether the two overlap.

\subsection{Results}
\ExpResults

\textbf{Alone on the champion} the knob helps a little: every value raises the panel finish score in the sweep (best 0.5: 79.1\,\% against 74.5\,\% on seven models), but on the unseen final-test maps the gain shrinks to 1.1 points on the fixed six (79.9\,\% against 78.8\,\%) and the head-to-head is 54.7\,\% finish, 55\,\% won ($\pm$5), which does not clear 50\,\%.

\textbf{On top of movement noise} it adds nothing: no value beats the base row in the sweep, and the final test of 0.5 gives 48.8\,\% finish and 49\,\% won ($\pm$5) against the base, with 81.6\,\% against 81.3\,\% on the fixed six.

\subsection{Behaviour}
The six stalled 8-player games of the diagnosis, replayed with \knob{spread} $=0.5$ alone: four now finish (turns 47 to 59), none did before. In the two that still stall, the leader's army is as jammed as before (58\,\% of its units cannot move or attack; 68\,\% have six or more own neighbours).

\todo{jakob}{What you saw in the live games, if you watched any.}

\subsection{Interpretation}
Spread and movement noise act on the same jam and largely do the same job: both let a boxed-in unit take a step that gains nothing toward its goal. Once noise does that, the extra pull away from own units has nothing left to fix, and at higher values it costs points (units leave useful positions). Alone, spread is the weaker of the two: 54.7\,\% against 60.1\,\% finish head to head.

\todo{jakob}{Your reading: agree, disagree, what is missing?}

\subsection{Threats to validity}
\begin{itemize}
\item Only the 4-player setup was swept; crowding is worst in 8-player games, where spread might add more.
\item Crowds are counted within one tile only; a wider radius or a cap per region (the formation idea, \issue{132}) is a different mechanism and untested.
\item The sweep picks its best point on 200 games per opponent; the gap between sweep and final test (79.1 against 76.3\,\% on seven models) shows how much of a sweep's best value is selection luck.
\end{itemize}

\subsection{Verdict and next step}
\textbf{Not kept.} Alone it fails both keep rules narrowly (head-to-head interval touches 50\,\%; panel gain 1.1 points, below 1.5). On top of movement noise it is a tie. The code stays, switched off. Next: formations with a strength cap (\issue{132}), which is the structured version of this idea.

\subsection{Notes for the paper}
Draft: two mechanisms aimed at the same failure are substitutes, not complements; measured one at a time on top of each other, the second adds nothing. A caution against crediting each of several simultaneous changes with the joint gain.

\todo{jakob}{Is this worth a place in the paper? Where?}


\clearpage\phantomsection\addcontentsline{toc}{section}{Experiment 9: Saving stars for a better purchase (not kept)}

% generated by scripts/experiment-report.ts from experiment.json and models/<run>/ - do not edit
\providecommand{\ExpHeader}{}\renewcommand{\ExpHeader}{%
{\sffamily\footnotesize\color{expgrey}EXPERIMENT 9 \textperiodcentered\ MECHANISM \hfill 2026-11-24}\par
{\LARGE\bfseries\sffamily Saving stars for a better purchase}\hfill\verdictbox{expred}{NOT KEPT}\par\smallskip
{\itshape Does a one-step greedy bot gain from skipping a cheap purchase that would leave it unable to afford a better one next turn?}\par\medskip
\begin{tabular}{@{}p{22mm}p{\dimexpr\linewidth-26mm}@{}}
\fact{Change}{A purchase that leaves next turn's stars short of the best purchase then affordable, and scores below it, loses reserve times that purchase's score.}
\fact{Knobs}{\knob{reserve} $\in$ \{0.15, 0.3, 0.6, 1\}, best 0.15}
\fact{Base}{\model{greedy-full-front}}
\fact{Method}{Strength sweep, 200 games per point against each panel model, final test 400 games; metric: points share}
\fact{Setups}{4 players, 16$\times$16, drylands, turn cap 150}
\fact{Record}{issue \issue{41}, pull request \issue{126}; status: done}
\end{tabular}\par\medskip
\begin{center}\begin{tikzpicture}[node distance=5mm]
    \node[flowbox] (n0) {{\sffamily\bfseries\scriptsize\color{expgrey}BASE}\\[1pt]\model{full-front}};
    \node[flowbox,right=5mm of n0] (n1) {{\sffamily\bfseries\scriptsize\color{expgrey}CHANGE}\\[1pt]\knob{reserve}\\{\scriptsize 0.15, 0.3, 0.6, 1}};
    \node[flowbox,right=5mm of n1] (n2) {{\sffamily\bfseries\scriptsize\color{expgrey}MEASURED}\\[1pt]Sweep, 200 games per point\\{\scriptsize 1 setup, metric: points share}};
    \node[flowbox,right=5mm of n2] (n3) {{\sffamily\bfseries\scriptsize\color{expgrey}RESULT}\\[1pt]Points vs base, final test\\{\scriptsize 50.0 $\rightarrow$ 50.0\,\%}\\\textbf{+0.0\pp}};
    \node[flowbox,right=5mm of n3] (n4) {{\sffamily\bfseries\scriptsize\color{expgrey}VERDICT}\\[1pt]\textbf{\color{expred}NOT KEPT}};
    \draw[flowarrow] (n0) -- (n1); \draw[flowarrow] (n1) -- (n2); \draw[flowarrow] (n2) -- (n3); \draw[flowarrow] (n3) -- (n4);
\end{tikzpicture}\end{center}
\begin{center}\sffamily\small\begin{tabular}{@{}lrrr@{}}\toprule
Measure & Base & New & Change\\\midrule
Points vs base, final test & 50.0\,\% & 50.0\,\% & +0.0\\
Points share, 6-model panel, sweep & 70.3\,\% & 70.3\,\% & +0.0\\
Decisions changed per 1,000 & 0.0 & 2.0 & +2.0\\
\bottomrule\end{tabular}\end{center}
}
\providecommand{\ExpResults}{}\renewcommand{\ExpResults}{%
\subsection{Sweep, 4 players {\normalfont\small(\model{greedy-full-reserve})}}
\begin{center}\sffamily\small\begin{tabular}{@{}lrrrrrr@{}}\toprule
Point & Finish (\%) & Points (\%) & Won (\%) & Conquest (\%) & Turns & vs base: points (\%)\\\midrule
base & -- & 67.36 & 80.7 & -- & -- & 50.0\\
reserve=0.15 & -- & 67.40 & 80.8 & -- & -- & 50.0\\
reserve=0.3 & -- & 67.40 & 80.8 & -- & -- & 50.0\\
reserve=0.6 & -- & 67.36 & 80.9 & -- & -- & 48.8\\
reserve=1 & -- & 67.31 & 81.2 & -- & -- & 49.1\\
\bottomrule\end{tabular}\\[3pt]{\footnotesize\color{expgrey}Panel means; 200 games against each of 7 panel models per point.}\end{center}
\begin{center}\begin{tikzpicture}\begin{axis}[expplot,height=4.6cm,xtick={0,1,2,3,4},xticklabels={{base},{0.15},{0.3},{0.6},{1}},
  xlabel={\texttt{reserve}},ylabel={Points share, panel mean (\%)},enlarge x limits=0.12,enlarge y limits=0.25]
\addplot[dashed,expgrey,line width=0.7pt,forget plot] coordinates {(0,67.357) (4,67.357)};
\addplot[only marks,mark=*,mark size=2.6pt,expblue] coordinates { (1,67.400) (2,67.400) (3,67.357) (4,67.314)};
\addplot[only marks,mark=*,mark size=2.6pt,expgrey] coordinates {(0,67.357)};
\end{axis}\end{tikzpicture}\\[-2pt]{\footnotesize\color{expgrey}Dashed: the base. Grey: base; blue: sweep points.}\end{center}
\begin{center}\sffamily\small\begin{tabular}{@{}lrr@{}}\toprule
Opponent & Won (\%) & Points (\%)\\\midrule
\model{v1} & 94 $\pm$ 2 & 75.6\\
\model{hill-v1} & 91 $\pm$ 3 & 72.3\\
\model{hill-v2} & 95 $\pm$ 2 & 73.7\\
\model{hill-v3} & 91 $\pm$ 3 & 70.3\\
\model{hill-v4} & 87 $\pm$ 3 & 68.6\\
\model{full-v1} & 58 $\pm$ 5 & 54.8\\
\model{full-front} & 50 $\pm$ 5 & 50.0\\
\midrule
\textbf{Mean without \model{full-front}} & 85.9 & 69.2\\
\bottomrule\end{tabular}\\[3pt]{\footnotesize\color{expgrey}Final test of reserve=0.15: 400 unseen games against each.}\end{center}
}

\ExpHeader

\subsection{Question and hypothesis}
The greedy bot spends stars on the best affordable action every step. The Autonomous Systems course calls this kind of one-step choice myopic and motivates planning as looking ahead and delaying action (AS 2.0 Planning as Search, PDF page 15, slide 8/51; technique AS-05, issue \#41). \textbf{Hypothesis:} skipping a cheap purchase that would block a clearly better one next turn raises the points share against the panel by about one point and wins the head-to-head against the base.

\subsection{Mechanism}
Each turn the bot finds the best purchase it cannot afford now but can afford next turn (stars set to next turn's income, cost in $(\mathit{stars}, \mathit{next}]$), with score $R$ and cost $C$. A purchase with cost $c>0$ and score $0<s<R$ for which $\mathit{next}-c<C$ scores $s-\knob{reserve}\cdot R$ instead. Default 0 (off). Code: \texttt{purchase\_cost}, \texttt{reserve\_reference} and the saving branch in \texttt{full\_best\_of}, \texttt{engine/core/src/bot\_full.rs}.

\subsection{Setup}
Values were sized so that at 1.0 a blocked purchase loses its full reference score, enough to push any cheaper purchase below 0. This experiment ran before the finish metric (\#124), so it uses the points share. Base: \model{greedy-full-front}, the champion at the time.

\subsection{Results}
\ExpResults

All four points lie within 0.2 points of the base on the 6-model panel (70.25\,\%), inside the noise of about $\pm$1.2. The final test of the best point (0.15) against the base is an exact tie: 50.0\,\% points, 50\,\% won ($\pm$5).

\subsection{Behaviour}
At \knob{reserve} 1.0 the mechanism changed 4 of 1,879 bot decisions (about 2 per 1,000) in the first 12 turns of generated 4-player games. The knob saturates at 1.0: larger values change nothing more.

\subsection{Interpretation}
The prediction failed because the condition almost never holds: a better purchase affordable next turn but not now, blocked by a cheaper one now. The idea is sound, but a one-turn lookahead on a single reference purchase is too narrow to change play. A plan over several turns is the DP of experiment 10.

\todo{jakob}{Your reading: agree, disagree, what is missing?}

\subsection{Threats to validity}
\begin{itemize}
\item The behaviour count comes from the first 12 turns only; later turns, with more stars, may trigger it more often.
\item The reference is one purchase. A plan that weighs several future purchases might act often enough to show an effect.
\end{itemize}

\subsection{Verdict and next step}
\textbf{Not kept}: the head-to-head is a tie and the panel moves less than the noise. Merged off by default (PR \#126) as the record. Champion stays \model{greedy-full-front}. Next: the multi-turn version, an economy plan by dynamic programming (experiment 10, \#68).

\subsection{Notes for the paper}
Draft: a negative result worth one line. A myopia patch that looks one turn ahead at one purchase almost never fires (2 per 1,000 decisions); the gain needs a real multi-step plan (experiment 10).

\todo{jakob}{Is this worth a place in the paper? Where?}


\clearpage\phantomsection\addcontentsline{toc}{section}{Experiment 10: Economy plan by dynamic programming (kept)}

% generated by scripts/experiment-report.ts from experiment.json and models/<run>/ - do not edit
\providecommand{\ExpHeader}{}\renewcommand{\ExpHeader}{%
{\sffamily\footnotesize\color{expgrey}EXPERIMENT 10 \textperiodcentered\ MECHANISM \hfill 2026-11-24}\par
{\LARGE\bfseries\sffamily Economy plan by dynamic programming}\hfill\verdictbox{expgreen}{KEPT}\par\smallskip
{\itshape Does scoring research and economy builds by a finite-horizon DP over one city's economy beat the greedy bot's fixed build and research scores?}\par\medskip
\begin{tabular}{@{}p{22mm}p{\dimexpr\linewidth-26mm}@{}}
\fact{Change}{Research, pop builds, leveling buildings and markets add plan times the DP advantage of buying now rather than next round, in stars, clamped to $\pm$10.}
\fact{Knobs}{\knob{plan} $\in$ \{1, 2, 4, 8\}, chosen 2}
\fact{Base}{\model{greedy-full-peacebreak} $\rightarrow$ result \model{greedy-full-economy}}
\fact{Method}{Strength sweep, 200 games per point against each panel model, final test 400 games; metric: finish score}
\fact{Setups}{4 players, 16$\times$16, drylands, turn cap 150}
\fact{Record}{issue \issue{68}, pull request \issue{135}; status: done}
\end{tabular}\par\medskip
\begin{center}\begin{tikzpicture}[node distance=5mm]
    \node[flowbox] (n0) {{\sffamily\bfseries\scriptsize\color{expgrey}BASE}\\[1pt]\model{full-peacebreak}};
    \node[flowbox,right=5mm of n0] (n1) {{\sffamily\bfseries\scriptsize\color{expgrey}CHANGE}\\[1pt]\knob{plan}\\{\scriptsize 1, 2, 4, 8}};
    \node[flowbox,right=5mm of n1] (n2) {{\sffamily\bfseries\scriptsize\color{expgrey}MEASURED}\\[1pt]Sweep, 200 games per point\\{\scriptsize 1 setup, metric: finish score}};
    \node[flowbox,right=5mm of n2] (n3) {{\sffamily\bfseries\scriptsize\color{expgrey}RESULT}\\[1pt]Finish score, 6-model panel, final-test maps\\{\scriptsize 78.8 $\rightarrow$ 82.5\,\%}\\\textbf{+3.7\pp}};
    \node[flowbox,right=5mm of n3] (n4) {{\sffamily\bfseries\scriptsize\color{expgrey}VERDICT}\\[1pt]\textbf{\color{expgreen}KEPT}\\{\scriptsize new champion:}\\\model{full-economy}};
    \draw[flowarrow] (n0) -- (n1); \draw[flowarrow] (n1) -- (n2); \draw[flowarrow] (n2) -- (n3); \draw[flowarrow] (n3) -- (n4);
\end{tikzpicture}\end{center}
\begin{center}\sffamily\small\begin{tabular}{@{}lrrr@{}}\toprule
Measure & Base & New & Change\\\midrule
Finish score, 6-model panel, final-test maps & 78.8\,\% & 82.5\,\% & +3.7\\
Finish vs base, final test & 50.0\,\% & 57.0\,\% & +7.0\\
Points vs base, final test & 50.0\,\% & 54.7\,\% & +4.7\\
Conquest wins, 6-model panel & 76.5\,\% & 81.0\,\% & +4.5\\
Game length in turns & 62.9 & 59.7 & $-$3.2\\
\bottomrule\end{tabular}\end{center}
}
\providecommand{\ExpResults}{}\renewcommand{\ExpResults}{%
\subsection{Sweep, 4 players {\normalfont\small(\model{greedy-full-economy})}}
\begin{center}\sffamily\small\begin{tabular}{@{}lrrrrrr@{}}\toprule
Point & Finish (\%) & Points (\%) & Won (\%) & Conquest (\%) & Turns & vs base: points (\%)\\\midrule
base & 74.46 & 67.93 & 80.6 & 72.5 & 60.5 & 50.0\\
plan=1 & 78.99 & 70.61 & 85.8 & 77.4 & 59.6 & 55.1\\
plan=2 & 79.26 & 71.26 & 87.1 & 77.3 & 60.1 & 54.6\\
plan=4 & 78.91 & 70.96 & 86.4 & 76.5 & 60.2 & 53.0\\
plan=8 & 78.10 & 69.96 & 84.7 & 76.4 & 58.2 & 49.5\\
\bottomrule\end{tabular}\\[3pt]{\footnotesize\color{expgrey}Panel means; 200 games against each of 7 panel models per point.}\end{center}
\begin{center}\begin{tikzpicture}\begin{axis}[expplot,height=4.6cm,xtick={0,1,2,3,4},xticklabels={{base},{1},{2},{4},{8}},
  xlabel={\texttt{plan}},ylabel={Finish score, panel mean (\%)},enlarge x limits=0.12,enlarge y limits=0.25]
\addplot[dashed,expgrey,line width=0.7pt,forget plot] coordinates {(0,74.457) (4,74.457)};
\addplot[only marks,mark=*,mark size=2.6pt,expblue] coordinates { (1,78.986) (2,79.257) (3,78.914) (4,78.100)};
\addplot[only marks,mark=*,mark size=2.6pt,expgrey] coordinates {(0,74.457)};
\end{axis}\end{tikzpicture}\\[-2pt]{\footnotesize\color{expgrey}Dashed: the base. Grey: base; blue: sweep points.}\end{center}
\begin{center}\sffamily\small\begin{tabular}{@{}lrrrrr@{}}\toprule
Opponent & Won (\%) & Points (\%) & Finish (\%) & Conquest (\%) & Turns\\\midrule
\model{v1} & 99 $\pm$ 1 & 79.9 & 90.0 & 91 & 52\\
\model{hill-v1} & 96 $\pm$ 2 & 76.2 & 86.7 & 88 & 56\\
\model{hill-v2} & 96 $\pm$ 2 & 76.1 & 86.1 & 85 & 61\\
\model{hill-v3} & 95 $\pm$ 2 & 73.5 & 85.5 & 86 & 58\\
\model{hill-v4} & 96 $\pm$ 2 & 73.9 & 83.7 & 79 & 66\\
\model{full-v1} & 66 $\pm$ 5 & 59.8 & 63.1 & 58 & 64\\
\model{full-peacebreak} & 58 $\pm$ 5 & 54.7 & 57.0 & 53 & 59\\
\midrule
\textbf{Mean without \model{full-peacebreak}} & 91.1 & 73.2 & 82.5 & 81.0 & 59.7\\
\bottomrule\end{tabular}\\[3pt]{\footnotesize\color{expgrey}Final test of plan=2: 400 unseen games against each.}\end{center}

\subsection{The base on the final-test maps, 4 players {\normalfont\small(\model{greedy-full-peacebreak-final})}}
\begin{center}\sffamily\small\begin{tabular}{@{}lrrrrr@{}}\toprule
Opponent & Won (\%) & Points (\%) & Finish (\%) & Conquest (\%) & Turns\\\midrule
\model{v1} & 96 $\pm$ 2 & 77.9 & 87.7 & 90 & 54\\
\model{hill-v1} & 91 $\pm$ 3 & 73.6 & 82.2 & 80 & 63\\
\model{hill-v2} & 95 $\pm$ 2 & 75.0 & 84.8 & 83 & 65\\
\model{hill-v3} & 91 $\pm$ 3 & 72.3 & 83.4 & 84 & 61\\
\model{hill-v4} & 88 $\pm$ 3 & 69.8 & 77.9 & 71 & 72\\
\model{full-v1} & 57 $\pm$ 5 & 54.7 & 56.7 & 52 & 63\\
\midrule
\textbf{Mean} & 86.2 & 70.5 & 78.8 & 76.5 & 62.9\\
\bottomrule\end{tabular}\\[3pt]{\footnotesize\color{expgrey}Final test of base: 400 unseen games against each.}\end{center}
}

\ExpHeader

\subsection{Question and hypothesis}
The greedy bot scores research and economy builds with fixed weights, so it cannot tell a build that pays back in five rounds from one that never does. The Reinforcement Learning course solves finite-horizon problems by backward induction: value equals immediate reward plus the value of what follows (RL 2.0 Dynamic Programming, PDF page 63, ``Easy answer for finite-horizon problems''; issue \#68). \textbf{Hypothesis:} scoring purchases by a DP over the economy raises the finish score on the 6-model panel by more than 1.5 points and wins the head-to-head against the base.

\subsection{Mechanism}
\texttt{econ\_plan.rs} solves a finite-horizon DP once per process (about 5\,ms) over one representative city: rounds left, cumulative population, buildable population supply and stars per city. For a purchase with cost $c$ it reads the advantage
\[
A = (1-\gamma)\,\bigl(V(\text{buy now}) - V(\text{never buy}) - \lambda c\bigr),
\]
the value of buying now rather than next round, in stars, clamped to $\pm 10$. Research, population builds, leveling buildings and markets score $s + \knob{plan}\cdot A$; training and upgrades are unchanged. Default 0 (off), and at 0 the scores are bit-identical. Code: \texttt{advantage} in \texttt{engine/core/src/econ\_plan.rs}, called from \texttt{score\_full\_with} in \texttt{bot\_full.rs}.

\subsection{Setup}
Values were sized so that $\knob{plan}\cdot A$ at plan 2 can tip a build against a move (about 25 plus tiles gained). The base was measured again on the final-test maps (\model{greedy-full-peacebreak-final}), so the panel means of base and new model come from the same maps.

\subsection{Results}
\ExpResults

Every value beats the base on the panel: finish 78.5\,\% (base) against 82.7, 83.1, 82.9 and 82.8\,\% at plan 1, 2, 4 and 8. Head to head against the base the effect fades as plan grows (finish 57.0, 56.3, 54.8 and 50.1\,\%). Final test of plan 2, 400 games: against the base 57.0\,\% finish, 54.7\,\% points, 58\,\% won ($\pm$5). On the final-test maps the panel finish rises from 78.8 to 82.5\,\%, conquest wins from 76.5 to 81.0\,\%, and games are 3 turns shorter.

\subsection{Behaviour}
The plan changes about 21 of 1,000 decisions. Forges, sawmills and markets get built where the bot used to end its turn, and research follows payback instead of noise. Play is about 8\,\% slower.

\todo{jakob}{What you saw in the live games, if you watched any.}

\subsection{Interpretation}
The prediction held, by both keep rules. An economy that compounds rewards early spending that pays back within the game's horizon, and the DP prices that horizon. The fade at plan 8 suggests that a large weight lets the economy crowd out military play.

\todo{jakob}{Your reading: agree, disagree, what is missing?}

\subsection{Threats to validity}
\begin{itemize}
\item The DP models one representative city, with no war and no losses; its values are a proxy, not the true game's value.
\item Only the 4-player setup was swept; 8-player maps have longer games, where the horizon matters more.
\item $\gamma$ and $\lambda$ were set by hand, not swept.
\end{itemize}

\subsection{Verdict and next step}
\textbf{Kept} by both rules (head-to-head and panel). New champion: \model{greedy-full-economy}, PR \#135. Next: it was never measured against experiment 7's champion, \model{greedy-full-movenoise}, which was kept against the same base. Experiment 11 does that.

\subsection{Notes for the paper}
Draft: dynamic programming over an abstract economy, solved once in 5\,ms, adds 4 points of finish score to a hand-tuned greedy bot. A course method carries over almost unchanged to a real game. Candidate figure: the sweep curve of the finish score over plan.

\todo{jakob}{Is this worth a place in the paper? Where?}


\clearpage\phantomsection\addcontentsline{toc}{section}{Experiment 11: Combining the two champions: economy plan plus movement noise (kept)}

% generated by scripts/experiment-report.ts from experiment.json and models/<run>/ - do not edit
\providecommand{\ExpHeader}{}\renewcommand{\ExpHeader}{%
{\sffamily\footnotesize\color{expgrey}EXPERIMENT 11 \textperiodcentered\ MECHANISM \hfill 2026-11-25}\par
{\LARGE\bfseries\sffamily Combining the two champions: economy plan plus movement noise}\hfill\verdictbox{expgreen}{KEPT}\par\smallskip
{\itshape Experiments 7 and 10 were both kept against greedy-full-peacebreak and never met: which is stronger, and does movement noise still help on top of the economy plan?}\par\medskip
\begin{tabular}{@{}p{22mm}p{\dimexpr\linewidth-26mm}@{}}
\fact{Change}{greedy-full-economy (plan 2) with moveNoise added: a zero-mean random term uniform in [$-$moveNoise, +moveNoise] on every move.}
\fact{Knobs}{\knob{moveNoise} $\in$ \{2, 4, 8\}, chosen 2}
\fact{Base}{\model{greedy-full-economy} $\rightarrow$ result \model{greedy-full-combo}}
\fact{Method}{Strength sweep, 200 games per point against each panel model plus both champions, final test 400 games; metric: finish score}
\fact{Setups}{4 players, 16$\times$16, drylands, turn cap 150}
\fact{Record}{issue \issue{31}; status: done}
\end{tabular}\par\medskip
\begin{center}\begin{tikzpicture}[node distance=5mm]
    \node[flowbox] (n0) {{\sffamily\bfseries\scriptsize\color{expgrey}BASE}\\[1pt]\model{full-economy}};
    \node[flowbox,right=5mm of n0] (n1) {{\sffamily\bfseries\scriptsize\color{expgrey}CHANGE}\\[1pt]\knob{moveNoise}\\{\scriptsize 2, 4, 8}};
    \node[flowbox,right=5mm of n1] (n2) {{\sffamily\bfseries\scriptsize\color{expgrey}MEASURED}\\[1pt]Sweep, 200 games per point\\{\scriptsize 1 setup, metric: finish score}};
    \node[flowbox,right=5mm of n2] (n3) {{\sffamily\bfseries\scriptsize\color{expgrey}RESULT}\\[1pt]Finish score, 6-model panel, final-test maps\\{\scriptsize 82.5 $\rightarrow$ 83.8\,\%}\\\textbf{+1.3\pp}};
    \node[flowbox,right=5mm of n3] (n4) {{\sffamily\bfseries\scriptsize\color{expgrey}VERDICT}\\[1pt]\textbf{\color{expgreen}KEPT}\\{\scriptsize new champion:}\\\model{full-combo}};
    \draw[flowarrow] (n0) -- (n1); \draw[flowarrow] (n1) -- (n2); \draw[flowarrow] (n2) -- (n3); \draw[flowarrow] (n3) -- (n4);
\end{tikzpicture}\end{center}
\begin{center}\sffamily\small\begin{tabular}{@{}lrrr@{}}\toprule
Measure & Base & New & Change\\\midrule
Finish score, 6-model panel, final-test maps & 82.5\,\% & 83.8\,\% & +1.3\\
Finish vs greedy-full-economy, final test & 50.0\,\% & 55.5\,\% & +5.5\\
Finish vs greedy-full-movenoise, final test & 50.0\,\% & 55.7\,\% & +5.7\\
Points share, 6-model panel & 73.2\,\% & 75.1\,\% & +1.9\\
Game length in turns & 59.7 & 59.1 & $-$0.6\\
\bottomrule\end{tabular}\end{center}
}
\providecommand{\ExpResults}{}\renewcommand{\ExpResults}{%
\subsection{Sweep, 4 players {\normalfont\small(\model{greedy-full-combo})}}
\begin{center}\sffamily\small\begin{tabular}{@{}lrrrrrr@{}}\toprule
Point & Finish (\%) & Points (\%) & Won (\%) & Conquest (\%) & Turns & vs base: points (\%)\\\midrule
base & 74.55 & 67.70 & 81.0 & 71.1 & 61.8 & 50.0\\
moveNoise=2 & 77.34 & 70.24 & 84.8 & 74.1 & 59.1 & 55.6\\
moveNoise=4 & 76.89 & 69.77 & 82.8 & 75.1 & 58.1 & 51.7\\
moveNoise=8 & 76.59 & 69.39 & 82.3 & 75.8 & 57.7 & 50.6\\
\bottomrule\end{tabular}\\[3pt]{\footnotesize\color{expgrey}Panel means; 200 games against each of 8 panel models per point.}\end{center}
\begin{center}\begin{tikzpicture}\begin{axis}[expplot,height=4.6cm,xtick={0,1,2,3},xticklabels={{base},{2},{4},{8}},
  xlabel={\texttt{moveNoise}},ylabel={Finish score, panel mean (\%)},enlarge x limits=0.12,enlarge y limits=0.25]
\addplot[dashed,expgrey,line width=0.7pt,forget plot] coordinates {(0,74.550) (3,74.550)};
\addplot[only marks,mark=*,mark size=2.6pt,expblue] coordinates { (1,77.337) (2,76.888) (3,76.587)};
\addplot[only marks,mark=*,mark size=2.6pt,expgrey] coordinates {(0,74.550)};
\end{axis}\end{tikzpicture}\\[-2pt]{\footnotesize\color{expgrey}Dashed: the base. Grey: base; blue: sweep points.}\end{center}
\begin{center}\sffamily\small\begin{tabular}{@{}lrrrrr@{}}\toprule
Opponent & Won (\%) & Points (\%) & Finish (\%) & Conquest (\%) & Turns\\\midrule
\model{v1} & 98 $\pm$ 2 & 80.5 & 88.3 & 87 & 54\\
\model{hill-v1} & 99 $\pm$ 1 & 79.4 & 88.9 & 89 & 54\\
\model{hill-v2} & 97 $\pm$ 2 & 77.2 & 87.5 & 86 & 57\\
\model{hill-v3} & 95 $\pm$ 2 & 74.4 & 85.2 & 83 & 60\\
\model{hill-v4} & 97 $\pm$ 2 & 75.1 & 84.1 & 77 & 68\\
\model{full-v1} & 73 $\pm$ 4 & 63.9 & 68.9 & 65 & 61\\
\model{full-economy} & 55 $\pm$ 5 & 53.9 & 55.5 & 50 & 63\\
\model{full-movenoise} & 57 $\pm$ 5 & 54.3 & 55.7 & 51 & 59\\
\midrule
\textbf{Mean without \model{full-economy}} & 87.8 & 72.1 & 79.8 & 76.8 & 59.0\\
\bottomrule\end{tabular}\\[3pt]{\footnotesize\color{expgrey}Final test of moveNoise=2: 400 unseen games against each.}\end{center}

\subsection{Base (economy) on the final-test maps {\normalfont\small(\model{greedy-full-economy-final})}}
\begin{center}\sffamily\small\begin{tabular}{@{}lrrrrr@{}}\toprule
Opponent & Won (\%) & Points (\%) & Finish (\%) & Conquest (\%) & Turns\\\midrule
\model{v1} & 99 $\pm$ 1 & 79.9 & 90.0 & 91 & 52\\
\model{hill-v1} & 96 $\pm$ 2 & 76.2 & 86.7 & 88 & 56\\
\model{hill-v2} & 96 $\pm$ 2 & 76.1 & 86.1 & 85 & 61\\
\model{hill-v3} & 95 $\pm$ 2 & 73.5 & 85.5 & 86 & 58\\
\model{hill-v4} & 96 $\pm$ 2 & 73.9 & 83.7 & 79 & 66\\
\model{full-v1} & 66 $\pm$ 5 & 59.8 & 63.1 & 58 & 64\\
\midrule
\textbf{Mean} & 91.1 & 73.2 & 82.5 & 81.0 & 59.7\\
\bottomrule\end{tabular}\\[3pt]{\footnotesize\color{expgrey}Final test of base: 400 unseen games against each.}\end{center}

\subsection{greedy-full-movenoise on the final-test maps {\normalfont\small(\model{greedy-full-movenoise-final})}}
\begin{center}\sffamily\small\begin{tabular}{@{}lrrrrr@{}}\toprule
Opponent & Won (\%) & Points (\%) & Finish (\%) & Conquest (\%) & Turns\\\midrule
\model{v1} & 97 $\pm$ 2 & 79.5 & 88.8 & 92 & 49\\
\model{hill-v1} & 91 $\pm$ 3 & 75.3 & 84.0 & 84 & 58\\
\model{hill-v2} & 93 $\pm$ 3 & 75.0 & 85.2 & 85 & 59\\
\model{hill-v3} & 90 $\pm$ 3 & 72.3 & 82.7 & 82 & 60\\
\model{hill-v4} & 94 $\pm$ 2 & 73.8 & 82.9 & 78 & 65\\
\model{full-v1} & 67 $\pm$ 5 & 59.9 & 64.1 & 62 & 60\\
\midrule
\textbf{Mean} & 88.6 & 72.6 & 81.3 & 80.5 & 58.6\\
\bottomrule\end{tabular}\\[3pt]{\footnotesize\color{expgrey}Final test of base: 400 unseen games against each.}\end{center}
}

\ExpHeader

\subsection{Question and hypothesis}
Two sessions each produced a champion from \model{greedy-full-peacebreak} on the same evening: experiment 7 (movement noise, \model{greedy-full-movenoise}) and experiment 10 (economy plan, \model{greedy-full-economy}). They never played each other, so the chain of experiments had split in two. The two mechanisms act on different decisions (moves, purchases), so they should add up. \textbf{Hypothesis:} economy plus noise beats each champion head to head and adds more than 1 point of panel finish over the better one.

\subsection{Mechanism}
No new code: the branch of experiment 7 is merged with main, so all knobs are present. The model is \model{greedy-full-economy} (\knob{plan} 2) with \knob{moveNoise} added: every move gets a zero-mean random term, uniform in $[-\knob{moveNoise}, +\knob{moveNoise}]$, and a move that gains nothing towards the goal is no longer ruled out (see experiment 7).

\subsection{Setup}
The panel is the fixed six plus both champions, so each point is measured against both. The sweep's base row is \model{greedy-full-economy} unchanged, which gives the champions' direct comparison. Values 2, 4 and 8 cover the flat range found in experiment 7, where 4 was chosen. Both champions were measured again on the final-test maps (the extra runs below).

\subsection{Results}
\ExpResults

\textbf{The two champions against each other} (200 games, sweep maps): \model{greedy-full-economy} scores 47.9\,\% finish and 47.4\,\% points against \model{greedy-full-movenoise}, a small edge for movenoise inside the noise. On the 6-model panel the economy bot is ahead (final-test maps: 82.5 against 81.3\,\%). In short, they are about level.

\textbf{The combination:} every noise value beats both. Unlike experiment 7, the best value is the smallest, 2, which scores 56.3\,\% finish against economy and 55.2\,\% against movenoise. Final test of moveNoise 2, 400 games: against economy 55.5\,\% finish, 53.9\,\% points, 55\,\% won ($\pm$5); against movenoise 55.7\,\% finish, 54.3\,\% points, 57\,\% won ($\pm$5). On the final-test maps the 6-model panel finish is 83.8\,\%, against 82.5 for economy (+1.3) and 81.3 for movenoise (+2.5).

\subsection{Behaviour}
On the final-test maps against the panel, conquest wins are 81.1\,\% (economy 81.0, movenoise 80.5) and games last 59.1 turns (59.7 and 58.6). A behaviour count for unjammed units was not measured here; see experiment 7 for the diagnosis.

\subsection{Interpretation}
The mechanisms add up only in part: the combination beats each parent by about 5 points of finish head to head, but its panel gain over the economy bot (+1.3) is smaller than either mechanism's gain alone. Against the panel both mechanisms mostly fix the same thing, games that do not end, so they overlap there. The best noise is smaller than on the old base. A plausible reason is that a bot that builds more already has more ways to make progress, so less randomness is needed to break jams.

\todo{jakob}{Your reading: agree, disagree, what is missing?}

\subsection{Threats to validity}
\begin{itemize}
\item The won rate against economy is 55\,\% with an interval of about $\pm$4.9: the lower bound clears 50 only just.
\item Only the 4-player setup; experiment 7 also used 8 players, where noise helped about as much.
\item The champions' direct comparison is from 200 games on the sweep maps, not a 400-game final test.
\end{itemize}

\subsection{Verdict and next step}
\textbf{Kept.} Against \model{greedy-full-economy} by the head-to-head rule (finish and points above 50\,\%, won interval clear of 50, panel mean not lower). Against \model{greedy-full-movenoise} by both rules (panel +2.5). New single champion: \model{greedy-full-combo} (plan 2, moveNoise 2), which unifies the two lines. Next: a tuning run over all knobs from this champion (experiment 12), then the value-iteration field over tiles (\#76).

\subsection{Notes for the paper}
Draft: a methods point. Parallel experiments from a shared base must be reconciled by playing the resulting champions against each other and against their combination. Here the two mechanisms combined beat both parents, with diminishing returns on the panel.

\todo{jakob}{Is this worth a place in the paper? Where?}


\clearpage\phantomsection\addcontentsline{toc}{section}{Experiment 12: Perturbation probe: tune all knobs with a significance test (kept)}

% generated by scripts/experiment-report.ts from experiment.json and models/<run>/ - do not edit
\providecommand{\ExpHeader}{}\renewcommand{\ExpHeader}{%
{\sffamily\footnotesize\color{expgrey}EXPERIMENT 12 \textperiodcentered\ TUNING \hfill 2026-11-25}\par
{\LARGE\bfseries\sffamily Perturbation probe: tune all knobs with a significance test}\hfill\verdictbox{exporange}{RUNNING}\par\smallskip
{\itshape Can nudging every knob at random in every game, and moving only the knobs whose up and down games differ significantly, find gains that hill climbing cannot see?}\par\medskip
\begin{tabular}{@{}p{22mm}p{\dimexpr\linewidth-26mm}@{}}
\fact{Change}{A new trainer. Each sample nudges all knobs up or down by 25\% at random and plays the two opposite variations against each other and against a league; a knob moves half a nudge when its up-versus-down difference reaches z = 3.}
\fact{Knobs}{\knob{all non-zero knobs (about 200)}}
\fact{Base}{\model{greedy-full-combo}}
\fact{Method}{Tuning run of 100,000 games (10 rounds), against a hill-climbing control of the same budget; final test 400 games; metric: finish score}
\fact{Setups}{4 players, 16$\times$16, drylands, turn cap 150}
\fact{Record}{issue \issue{89}; status: running}
\end{tabular}\par\medskip
\begin{center}\begin{tikzpicture}[node distance=5mm]
    \node[flowbox] (n0) {{\sffamily\bfseries\scriptsize\color{expgrey}BASE}\\[1pt]\model{full-combo}};
    \node[flowbox,right=5mm of n0] (n1) {{\sffamily\bfseries\scriptsize\color{expgrey}CHANGE}\\[1pt]\knob{all non-zero knobs (about 200)}};
    \node[flowbox,right=5mm of n1] (n2) {{\sffamily\bfseries\scriptsize\color{expgrey}MEASURED}\\[1pt]100,000-game run vs a hill-climbing control\\{\scriptsize 1 setup, metric: finish score}};
    \node[flowbox,right=5mm of n2] (n3) {{\sffamily\bfseries\scriptsize\color{expgrey}RESULT}\\[1pt]--};
    \node[flowbox,right=5mm of n3] (n4) {{\sffamily\bfseries\scriptsize\color{expgrey}VERDICT}\\[1pt]\textbf{\color{exporange}RUNNING}};
    \draw[flowarrow] (n0) -- (n1); \draw[flowarrow] (n1) -- (n2); \draw[flowarrow] (n2) -- (n3); \draw[flowarrow] (n3) -- (n4);
\end{tikzpicture}\end{center}

}
\providecommand{\ExpResults}{}\renewcommand{\ExpResults}{%
\subsection{Perturbation probe, 100,000 games {\normalfont\small(\model{greedy-full-probe})}}
\begin{center}\begin{tikzpicture}\begin{axis}[expplot,height=4.6cm,xlabel={training games (thousands)},ylabel={Finish score, panel mean (\%)},enlarge y limits=0.2]
\addplot[dashed,expgrey,line width=0.7pt,forget plot] coordinates {(0.0,79.60) (0.0,79.60)};
\addplot[expblue,mark=*,mark size=1.8pt] coordinates {(0.0,79.60)};
\end{axis}\end{tikzpicture}\\[-2pt]{\footnotesize\color{expgrey}Learning curve of the tuning run (0 kept changes in 0 rounds). Dashed: the start.}\end{center}

\begin{center}\begin{tikzpicture}\begin{axis}[expplot,height=4.8cm,xmin=-5,xmax=5,ymin=0,xlabel={z-score of a knob (up games minus down games, in standard errors)},ylabel={knobs}]
\addplot[ybar interval,fill=expblue!55,draw=white] coordinates {(-5.0,3) (-4.5,0) (-4.0,0) (-3.5,2) (-3.0,4) (-2.5,5) (-2.0,13) (-1.5,19) (-1.0,24) (-0.5,37) (0.0,33) (0.5,28) (1.0,20) (1.5,8) (2.0,1) (2.5,1) (3.0,1) (3.5,0) (4.0,0) (4.5,6) (5.0,0)};
\addplot[dashed,expgrey,line width=0.8pt,domain=-5:5,samples=80] {205*0.5*exp(-x^2/2)/sqrt(2*pi)};
\draw[expgrey,line width=0.4pt] (axis cs:-3,0) -- (axis cs:-3,\pgfkeysvalueof{/pgfplots/ymax}) (axis cs:3,0) -- (axis cs:3,\pgfkeysvalueof{/pgfplots/ymax});
\end{axis}\end{tikzpicture}\\[-2pt]{\footnotesize\color{expgrey}z-scores of all 205 knobs after round 1 (10k training games). Dashed: pure chance. Vertical lines: the threshold $|z|=3$; 12 beyond it.}\end{center}
\begin{center}\sffamily\small\begin{tabular}{@{}rlrrr@{}}\toprule
Round & Knob & From & To & z\\\midrule
1 & \texttt{greedy.research} & 27.6 & 31.275 & $+7.6$\\
1 & \texttt{type.attack} & 1 & 1.133 & $+3.2$\\
1 & \texttt{tech.riding} & 0.948 & 0.837 & $-5.0$\\
1 & \texttt{tech.strategy} & 0.942 & 0.831 & $-3.1$\\
1 & \texttt{tech.climbing} & 1 & 1.133 & $+5.9$\\
1 & \texttt{tech.meditation} & 1.27 & 1.439 & $+6.6$\\
1 & \texttt{tech.philosophy} & 2.045 & 2.317 & $+5.8$\\
1 & \texttt{unit.warrior} & 1 & 1.133 & $+4.7$\\
1 & \texttt{unit.rider} & 1 & 0.882 & $-6.8$\\
1 & \texttt{reward.park} & 1 & 0.882 & $-7.3$\\
1 & \texttt{reward.superUnit} & 1 & 1.133 & $+7.6$\\
1 & \texttt{centre} & 1 & 0.882 & $-3.0$\\
\bottomrule\end{tabular}\\[3pt]{\footnotesize\color{expgrey}Knobs moved: 12.}\end{center}
\begin{center}\sffamily\small\begin{tabular}{@{}lrrrrrr@{}}\toprule
Knob & Value & Samples & Effect (pts) & z & z self & z league\\\midrule
\texttt{greedy.research} & 27.6 & 1667 & $+8.89$ & $+7.6$ & $+6.3$ & $+5.9$\\
\texttt{reward.superUnit} & 1 & 1667 & $+8.89$ & $+7.6$ & $+8.1$ & $+4.4$\\
\texttt{reward.park} & 1 & 1667 & $-8.59$ & $-7.3$ & $-8.2$ & $-3.9$\\
\texttt{unit.rider} & 1 & 1667 & $-7.96$ & $-6.8$ & $-6.3$ & $-4.7$\\
\texttt{tech.meditation} & 1.27 & 1667 & $+7.76$ & $+6.6$ & $+6.0$ & $+4.7$\\
\texttt{tech.climbing} & 1 & 1667 & $+6.89$ & $+5.9$ & $+5.5$ & $+4.0$\\
\texttt{tech.philosophy} & 2.045 & 1667 & $+6.81$ & $+5.8$ & $+4.5$ & $+4.8$\\
\texttt{tech.riding} & 0.948 & 1667 & $-5.84$ & $-5.0$ & $-4.4$ & $-3.6$\\
\texttt{unit.warrior} & 1 & 1667 & $+5.51$ & $+4.7$ & $+3.6$ & $+3.9$\\
\texttt{type.attack} & 1 & 1667 & $+3.82$ & $+3.2$ & $+3.4$ & $+1.9$\\
\bottomrule\end{tabular}\\[3pt]{\footnotesize\color{expgrey}Strongest evidence at the end of the run. Effect: finish score gained by one nudge up against one nudge down, halved.}\end{center}

\subsection{Control: hill climbing, 100,000 games {\normalfont\small(\model{greedy-full-probe-ctrl})}}
{\small\color{expgrey}No data recorded for this run yet.}
}

\ExpHeader

\subsection{Question and hypothesis}
Jakob, 25 November: ``Could we also add noise to other weights and then analyse the distribution and if we find a statistically significant difference between the positive and negative parts of the noise then we adjust the weights in that direction?''

Hill climbing accepts or rejects a whole bundle of changes on 300 to 400 games and found almost nothing in three runs of 100,000 games (experiments 1 to 3). The evolution-strategies trainer moves every weight on every step, significant or not. \textbf{Hypothesis:} measuring every knob in the same games and moving only the significant ones finds real single-knob gains that the other two trainers miss.

\textbf{Prediction, written before the run.} (1) Between 2 and 8 knobs move in 100,000 games, mostly research and tech weights (a 1,200-game smoke test already showed \knob{greedy.research} and three tech weights above $|z|=3$). (2) The panel finish score rises by 0.5 to 1.5 points. (3) Head to head against the base the result is 50 to 54\,\%, probably not clear of 50\,\% at $\pm$5. \textbf{Counts against the idea:} the z-scores of all knobs look like pure noise (spread near 1, nothing repeats after a knob's evidence is reset), or the panel curve falls as knobs move.

\subsection{Mechanism}
Each sample draws a sign $s_k\in\{-1,+1\}$ for every knob $k$, independently, and builds two opposite variations of the centre $\theta$: $w^{\pm}_k=\theta_k\,e^{\pm\sigma s_k}$ with $\sigma=0.25$ (a knob at 0 stays 0 and is not probed). On one map, every game from both seats, $w^+$ plays $w^-$, and both play the same league opponent. The outcome is
\[ x = \underbrace{\mathrm{finish}(w^+\text{ vs }w^-)-\tfrac12}_{\text{self}} \;+\; \underbrace{\mathrm{finish}(w^+\text{ vs }o)-\mathrm{finish}(w^-\text{ vs }o)}_{\text{league}} . \]
Because the signs are independent, the least-squares slope of $x$ on $s_k$ is $\hat\beta_k=\frac1n\sum_i x_i s_{ik}$: half the difference between the games where the knob was nudged up and those where it was nudged down. Its standard error is $\mathrm{sd}(x s_k)/\sqrt n$ and $z_k=\hat\beta_k/\mathrm{SE}$. After each round a knob with $|z_k|\ge 3$ moves half a nudge in the better direction and its evidence is reset; all other knobs keep accumulating evidence. Code: \texttt{scripts/train-probe.ts}. Course link: least squares and its analytic solution [EMAI ML-13, ML 2.0 p.27].

\subsection{Setup}
Base: the combined champion (economy plan plus movement noise 2), which beat both of its parents. 10 rounds of 1,667 samples (6 games each), league \model{hill-v4}, \model{full-v1} and the base. A curve point (200 games against each of the seven panel models, the same maps as every sweep) every two rounds; the centre of the best curve point is saved and plays the final test.

Sizing: one sample's outcome has a spread of about 48 points (smoke test), so after 100,000 games a knob that never moved has a standard error of about 0.4 points and is detected from about 1.1 points. With about 200 knobs and $|z|\ge 3$, about half a knob per round crosses the threshold by chance; such a knob has no real effect, so moving it should cost little, and the final test on unseen maps is the guard.

Control: hill climbing (\texttt{train-full}) from the same base with the skill's settings and the same 100,000 games.

\subsection{Results}
\ExpResults

\textbf{The probe.} The finish score against the seven panel models rose from 79.6\,\% to 82.5\,\% after two rounds and stayed between 82.3 and 83.0 from then on; the centre after round 8 (83.0\,\%) was saved. In the final test it won 61.3\,\% $\pm$4.8 of 400 games against the base (finish 59.0\,\%, points 56.2\,\%), and against the six fixed panel models its finish score is 85.80\,\% against the base's 83.82\,\%, with more conquests (85.5\,\% against 81.1\,\%) in shorter games (52.9 turns against 59.1).

\textbf{The moves.} 26 of 205 knobs moved, in 60 steps; 24 of the steps came in the first two rounds. The saved model holds the 52 steps up to round 8. The run is not noise: (i) in every one of the 60 moves the self-play comparison and the league comparison, which are different games, agree on the direction; (ii) \knob{defendRadius}, which cannot matter because \knob{defend} is 0, ended at $z=+0.07$; (iii) the 179 knobs that never moved have z-scores with a spread of 1.14 (chance: 1.00), 14 of them beyond $|z|=2$ where chance gives 8; (iv) only one knob, \knob{unit.warrior}, ever moved in both directions.

\textbf{The control.} Hill climbing with the same 100,000 games kept one challenger out of 297 (four small changes, panel finish 79.6\,\% to 79.7\,\%). Its saved model is the base itself: 50.0\,\% against the base, 83.82\,\% against the panel.

\textbf{8 players} (20$\times$20, final test only). The tuned model won 72.8\,\% $\pm$4.4 against the base (finish 67.9\,\%, points 64.6\,\%); against the six panel models its finish score is 87.13\,\% against the base's 83.08\,\%. The gain is larger than with 4 players, although no 8-player game was used in tuning.

\subsection{Behaviour}
\textbf{What changed in the bot.} Research is reweighted, not raised: the general research factor \knob{type.research} fell from 1 to 0.54 while the base research score \knob{greedy.research} rose from 27.6 to 35.4 and single technologies moved apart: philosophy 2.05 to 3.82, meditation 1.27 to 1.85, climbing 1 to 1.28; riding 0.95 to 0.65, fishing 1 to 0.69, strategy 0.94 to 0.73. Attacks count more (\knob{type.attack} 1 to 1.65). Among units, mind benders (2.00 to 2.56), warriors (1 to 1.28) and sirens went up, riders (1 to 0.78) and moonis down. Three mechanisms of earlier experiments moved slightly up or down: \knob{moveNoise} 2 to 2.27, \knob{centre} 1 to 0.78; \knob{front} moved only after the saved round.

\textbf{The sensitivity map.} After 100,000 games a knob that never moved has a standard error of 0.30 points, so an effect of 0.9 points per nudge would have been found. 179 of 205 knobs (87\,\%) are below that. The bot's strength, as far as a 25\,\% change can show, sits in about two dozen weights: research, a handful of technologies, unit choice and the attack factor. All building weights and almost all reward weights are in the insensitive group.

\textbf{A cliff, found and checked.} In round 3 \knob{unit.warrior} (then 1.284) showed $z=-22.3$ and \knob{unit.mindBender} (1.995) $z=+14.0$, and the spread of the outcome jumped from 48 to 58 points. The two effects have the size a pure interaction predicts: a loss only when the warrior weight is nudged up \emph{and} the mind bender weight down. The check run confirms it on the base (panel finish, seven models): warrior weight 1.649 alone gives 79.47\,\% (base 79.33\,\%), mind bender weight 1.554 alone plays exactly like the base (79.33\,\%), and both together fall to 54.21\,\% and take 30.4\,\% of the points against the base. So the bot collapses as soon as warriors outweigh mind benders in the unit choice. The trainer stepped back in the same round and the cliff did not recur once the mind bender weight had risen.

\todo{jakob}{What you saw in the live games, if you watched any.}

\subsection{Interpretation}
The prediction was wrong three times, each time too cautious: 26 knobs moved (predicted 2 to 8), the panel finish rose by 2.0 points (predicted 0.5 to 1.5), and the head-to-head result is 61\,\% (predicted 50 to 54\,\%). The prediction about \emph{where} the effects sit was right: research and technology weights lead.

Why the probe finds what hill climbing does not: hill climbing spends a whole round of 300 to 400 games on one question (is this bundle of random changes better?) and answers yes or no. The probe asks 205 questions of the same games, and because the signs are drawn independently, every knob's answer is an unbiased slope with the other 204 knobs acting as noise. That noise is the price: one sample's outcome spreads by 40 to 48 points, so only effects above about 3 points per nudge show within one round. The base evidently had many of those. The hand-set unit and technology weights of the base had never been tuned in this space, so the first tuning method able to see single knobs found a lot at once; a second run from the new champion should find much less.

The cliff shows the method's blind spot and its safety net at once: marginal effects cannot represent an interaction, but a step into a bad region produces an enormous z-score in the next round and is undone.

\todo{jakob}{Your reading: agree, disagree, what is missing?}

\subsection{Threats to validity}
\begin{itemize}
\item \textbf{One run, one seed.} The probe was run once. The final tests use unseen maps and the 8-player check a different format, which guards the strength claim, but which knobs move first may differ between seeds.
\item \textbf{The comparison of methods is not clean.} The control used the skill's hill-climbing settings (league: itself and its past), the probe a league of two panel models and the base. The probe therefore trained against opponents that are also in the test panel (\model{hill-v4}, \model{full-v1}); the head-to-head result against the base and the four other panel models do not have this overlap, and they improved as well.
\item \textbf{Multiple comparisons.} With 205 knobs and $|z|\ge 3$, about 0.55 false moves per round are expected, so roughly 5 of the 60 steps may be chance. They are half-nudges of knobs without a measurable effect; the final test bounds their cost.
\item \textbf{The saved round was chosen on the panel curve.} Picking the best of six curve points favours the trainer slightly on those maps; the final test uses other maps.
\item \textbf{The neighbourhood is noisy.} Effects are measured between variations that differ in all knobs by $\pm$25\,\%, not at the centre itself. A knob's effect at the centre can be smaller.
\end{itemize}

\subsection{Verdict and next step}
\textbf{Kept.} \model{greedy-full-probe} is the new champion: 61.3\,\% $\pm$4.8 won against the base (finish 59.0\,\%, points 56.2\,\%) and a panel finish of 85.80\,\% against 83.82\,\%; with 8 players 72.8\,\% $\pm$4.4. The hill-climbing control with the same budget found nothing.

The trainer stays as \texttt{scripts/train-probe.ts} (\texttt{npm run train-probe}). Next steps, each its own experiment: (1) a second probe run from the new champion with a smaller nudge, to see how much is left; (2) the same trainer on the plain greedy bots (issue \issue{131}); (3) the warrior and mind bender cliff as a bug report on the unit choice; (4) repeat with another seed to learn how stable the set of moved knobs is.

\subsection{Notes for the paper}
\textbf{Claim (draft).} Estimating all 205 weights' effects from the same games by random sign perturbation, and moving only the weights whose effect passes a significance test, improved a hand-built heuristic bot by 2.0 points of finish score and to 61\,\% against its predecessor in 100,000 games, where hill climbing with the same budget found nothing.

\textbf{Figures worth keeping.} The z-score histogram against the chance curve after round 1; the sensitivity map (87\,\% of weights without measurable effect); the cliff as a worked example of an interaction that marginal effects miss and the next round corrects. Course link: least squares on a designed experiment [EMAI ML-13, ML 2.0 p.27]; issue \issue{89}.

\todo{jakob}{Is this worth a place in the paper? Where?}


\end{document}
